\documentclass[11pt,letterpaper]{article}
\usepackage[T1]{fontenc}
\usepackage{lmodern,amsmath,amssymb,amsthm,mathtools}
\usepackage[margin=0.86in]{geometry}
\usepackage{booktabs,longtable,array,microtype,enumitem,xcolor}
\usepackage{fancyhdr,xurl,hyperref,listings}
\hypersetup{colorlinks=true,linkcolor=blue!45!black,citecolor=blue!45!black,urlcolor=blue!45!black,pdftitle={Local certificates for Diophantine consistency},pdfauthor={AI-assisted research for Vladimir Reshetnikov}}
\setlength{\parskip}{0.35em}
\setlength{\parindent}{0pt}
\emergencystretch=2em
\pagestyle{fancy}\fancyhf{}\fancyhead[L]{\small Local certificates for Diophantine consistency}\fancyhead[R]{\small ProveIt}\fancyfoot[C]{\thepage}
\setlength{\headheight}{14pt}
\newtheorem{theorem}{Theorem}[section]
\newtheorem{lemma}[theorem]{Lemma}
\newtheorem{proposition}[theorem]{Proposition}
\newtheorem{corollary}[theorem]{Corollary}
\theoremstyle{definition}\newtheorem{definition}[theorem]{Definition}
\theoremstyle{remark}\newtheorem{remark}[theorem]{Remark}
\newcommand{\Z}{\mathbb Z}\newcommand{\N}{\mathbb N}\newcommand{\R}{\mathbb R}
\newcommand{\Con}{\operatorname{Con}}\newcommand{\ZFC}{\mathsf{ZFC}}
\newcommand{\srcpin}{\path{1c5367452a0bd4dd7e9e4bceb1508b62dbd2ff61}}
\newcommand{\Ht}{\operatorname{ht}}\newcommand{\supp}{\operatorname{supp}}
\lstset{basicstyle=\ttfamily\small,breaklines=true,columns=fullflexible,keepspaces=true}
\newcommand{\DirectVars}{306}
\newcommand{\DirectLast}{305}
\newcommand{\DirectRows}{206}
\newcommand{\DirectRowLast}{205}
\newcommand{\DirectTerms}{6{,}259}
\newcommand{\DirectHeight}{128{,}446}

\title{Local certificates for Diophantine consistency\\[3pt]\large Shared Pell orientation and smaller quartics with unit coefficients}
\author{Research contribution prepared for Vladimir Reshetnikov\\\small Mathematical development, implementation, and review with AI assistance}
\date{24 October 2026 (America/Los\_Angeles)}
\begin{document}
\maketitle
\begin{abstract}
We give explicit ordinary integer polynomials whose lack of integer solutions
is provably equivalent in ZFC to the syntactic consistency of ZFC. The work
builds on the direct numerical-calculus branch of ProveIt and develops two
reductions. First, a linear trace-orientation term is added to an existing
pooled Pell-index congruence. This shares the work formerly performed by nine
separate sign certificates. Second, coefficient-preserving identifications of
graph aliases are certified by exact polynomial collection; a global
four-cycle restriction on an auxiliary factor graph is unnecessary. A
two-variable quartic equality energy removes further auxiliary variables.
The release contains the literal residuals and complete collected coefficient
lists, constructive zero-set arguments, deterministic reconstruction code,
and independent finite checkers. Every final variable ranges over all of
$\Z$. The reductions are ordinary mathematical proofs formalizable in ZFC,
composed with the explicitly identified inherited C222/ZF/ZFC calibration.
They do not claim a new end-to-end proof-assistant verification or a globally
simplest equation.
\end{abstract}

\section{The actual equations and the scope of the result}
\label{sec:outcome}

The principal direct equation is the fixed polynomial
\begin{equation}\label{eq:direct}
 P_{\mathrm{dir}}(z_0,\ldots,z_{\DirectLast})
   =\sum_{j=0}^{\DirectRowLast} r_j(z)^2=0,
 \qquad z\in\Z^{\DirectVars}.
\end{equation}
Here the $r_j$ are all the quadratic residuals in
\path{direct/integer306/equation.json}, in their specified coordinate order. Their
complete readable transcription is supplied with the release. Its exact
degree is four, its collected support has \DirectTerms\ monomials, and its
coefficient height is \DirectHeight. This is a finite literal equation:
there is no variable-exponent operation, binomial predicate, unspecified
universal polynomial, or additional implicit existential variable in
\eqref{eq:direct}.

The unit-coefficient equations have the particularly uniform form
\begin{equation}\label{eq:unit-literal}
 P_{\mathrm{unit}}(x)=1+
       \sum_{\alpha\in\mathcal A_+}x^\alpha
       -\sum_{\alpha\in\mathcal A_-}x^\alpha=0,
 \qquad x\in\Z^n.
\end{equation}
The two finite, disjoint sets of nonconstant exponent vectors are exactly
the positive and negative entries of each released
\path{polynomial.json}. Repeated coordinate indices mean powers; the empty
index tuple denotes the constant. The coefficient list itself, not a request
to generate some unspecified polynomial, defines the equation. The selected
profiles and their exact metrics are listed below.

\begin{table}[ht]
\centering\small
\begin{tabular}{@{}p{0.44\textwidth}rrrr@{}}\toprule
Construction & Variables & Degree & Terms & Height\\\midrule
Parent direct quartic [131] & 324 & 4 & 5{,}808 & 128{,}446\\
New direct quartic & 306 & 4 & 6{,}259 & 128{,}446\\
Parent unit-coefficient quartic [127] & 2{,}215 & 4 & 30{,}468 & 1\\
New graph, same 442-coordinate source & 1{,}252 & 4 & 27{,}738 & 1\\
New graph, fewer variables & 985 & 4 & 95{,}609 & 1\\
New graph, smaller support & 1{,}122 & 4 & 25{,}922 & 1\\
\bottomrule\end{tabular}
\caption{All variables are unrestricted integers. Terms are fully collected nonzero monomials. The graph companions permit individual fourth powers.}
\label{tab:metrics}\end{table}


For every selected polynomial $P$ of arity $n$, our conclusion is
\begin{equation}\label{eq:main}
 \ZFC\ \vdash\
 \bigl[\neg\exists x\in\Z^n\;P(x)=0\bigr]
       \ \longleftrightarrow\ \Con(\ZFC).
\end{equation}
The proof proceeds by proving the positive existential statement equivalent
to the existence of a finite syntactic ZFC refutation and then negating it.
No assertion of either side of \eqref{eq:main} is needed.

\subsection{What ``simpler'' means here}

There is no total ordering of all the relevant costs. We report arity,
degree, number of collected monomials, and coefficient height separately.
In particular, the repository already has quartics with fewer variables
and very large fixed coefficient tuples, and higher-degree equations with
far fewer variables \cite{repooverview}. The direct equation improves the
small-coefficient C222 construction; the graph equations improve the
unit-coefficient quartic construction. None is asserted to minimize arity
among all consistency polynomials.

For a nonzero polynomial with integer coefficients, coefficient height
cannot be less than one. Thus the graph equations have the smallest
possible nonzero coefficient height, while their arity and support remain
subjects for optimization. This elementary height observation is quite
different from a global simplicity theorem.

The source snapshot is \srcpin. The repository's main catalogue at that
snapshot still summarizes earlier waves, although packages 127--133 have
already been placed. We inspected those newer packages directly. Comparisons
are to specified literal artifacts, so a change of source, domain, or
individual-degree bound is stated explicitly.

\section{An elementary graph representation of equations}
\label{sec:graph}

We first describe the inherited graph representation, because it exposes
exactly what the new contraction theorem preserves. Let $\mathcal G=(V,E)$
be a finite simple graph, with an arbitrary orientation. Attach an integer
polynomial $m_e$ of degree at most two to each edge. Put
\[
 d_v=\sum_{e\text{ leaves }v}m_e-\sum_{e\text{ enters }v}m_e,
 \qquad \mathcal E_{\mathcal G}=\frac12\sum_{v\in V}d_v^2.
\]
The expression is nonnegative over $\R$ and vanishes precisely when every
vertex divergence vanishes. The factor $1/2$ is proof notation: after
collection the supplied final polynomial has integral coefficients.

Before substituting the edge polynomials, consider independent edge symbols
$t_e$. Each $t_e^2$ occurs at its two endpoints and therefore has coefficient
one in $\mathcal E_{\mathcal G}(t)$. Two distinct edges in a simple graph
share at most one endpoint, so their product has coefficient $+1$, $-1$, or
zero. The resulting height-one assertion need not survive an arbitrary
substitution $t_e=m_e$; different edge products may become the same
monomial. That is the only coefficient issue the new exact collector must
settle.

\begin{lemma}[Integral tree flows]\label{lem:tree}
Partition the vertices into connected blocks whose internal edges form
trees. Give all edges between blocks fixed integer values. There is an
integer assignment to the internal edges making every divergence zero if
and only if each block has total external signed flow zero. When it exists,
the assignment is unique.
\end{lemma}
\begin{proof}
Summing divergences over a block cancels all its internal edges, proving
necessity. For sufficiency, root its tree and remove leaves successively.
The only remaining internal edge at a leaf has coefficient $+1$ or $-1$,
so the leaf equation uniquely determines its integer value. Transfer this
contribution to the parent and continue. The final root equation is the
assumed block balance. The construction proves integrality and uniqueness.
\end{proof}

Normalize a source system into rows
\[
 R_r=\sum_{i=1}^{s_r}\epsilon_{ri}m_{ri},\qquad
       \epsilon_{ri}\in\{-1,1\},\quad\deg m_{ri}\leq2.
\]
The normalization uses monic topological equations to construct positive
integer coefficients by binary addition and to construct constant multiples
of source variables. The source coordinates are retained, every added
coordinate has a unique integer value, and the normalization therefore
preserves the integer zero set bijectively.

For a chosen positive integer $H$, use one hub block of $H$ vertices. The
block of row $r$ has $\lceil s_r/H\rceil$ vertices. Represent each signed
term occurrence by a distinct edge between its row block and the hub
block; the available bipartite capacity is sufficient. Connect each block
internally by a tree. Row-block balance is exactly $R_r=0$, and hub-block
balance is the negative sum of those equations. Lemma~\ref{lem:tree} gives
the desired equivalence. The number of independent tree-flow coordinates is
\begin{equation}\label{eq:flow-count}
 H-1+\sum_r\bigl(\lceil s_r/H\rceil-1\bigr).
\end{equation}

Each normalized logical variable is represented on its occurrences by one
or more \emph{aliases}. An alias is an ordinary independent integer
coordinate until its equality energy is imposed. Terms are homogenized to
degree two using aliases of the constant one. Tree-flow variables occupy a
separate coordinate class. There is no hidden interpretation imposed on
the quantified coordinates: all alias and unit conditions are equations
inside the final polynomial.

\section{A smaller equality certificate}
\label{sec:equality}

For $k\geq3$ aliases of one logical coordinate, use the cycle energy
\begin{equation}\label{eq:cycle}
 C(a_1,\ldots,a_k)=\sum_{i=1}^k a_i^2-
                      \sum_{i=1}^k a_i a_{i+1},\qquad a_{k+1}=a_1.
\end{equation}
Twice this polynomial is $\sum_i(a_i-a_{i+1})^2$. It is nonnegative and
vanishes exactly when all aliases agree. Its collected nonzero coefficients
are $+1$ or $-1$. A singleton needs no equality equation. A two-vertex
cycle would have coefficient $-2$ on $a_1a_2$, so the earlier compiler
inserted a third, dummy alias.

\begin{lemma}[Two-coordinate quartic equality]\label{lem:pair}
The polynomial
\begin{equation}\label{eq:pair}
 F(a,b)=a^4+b^4-a^3b-ab^3
\end{equation}
is nonnegative over $\R$ and vanishes exactly when $a=b$. Its four nonzero
coefficients have absolute value one. In particular it requires no dummy
coordinate.
\end{lemma}
\begin{proof}
Two useful exact identities are
\begin{align*}
 F(a,b)&=(a-b)^2(a^2+ab+b^2),\\
 2F(a,b)&=(a^2-b^2)^2+(a^2-ab)^2+(ab-b^2)^2.
\end{align*}
The first proves the equality condition because
$a^2+ab+b^2=(a+b/2)^2+3b^2/4$ is positive except at $(0,0)$, where
$a=b$ already holds. The second gives a direct rational
sum-of-squares certificate and is used by the independent collector.
\end{proof}

The new equality energy must be checked against the other quartic
monomials. For example, if $a^4$ already has coefficient $+1$ in the graph
energy, adding \eqref{eq:pair} would produce coefficient $+2$. Such a
replacement is rejected. The compiler applies the lemma only where the
complete collected polynomial keeps height one; otherwise it retains the
three-alias quadratic cycle. This is an exact finite condition, not an
assumption that the new monomials occupy a fresh namespace.

For unit aliases $u_1,\ldots,u_k$, with $k\geq2$, use
\begin{equation}\label{eq:unit-path}
 U=1+\sum_{i=1}^k u_i^2-u_1-u_k-
                          \sum_{i=1}^{k-1}u_i u_{i+1}.
\end{equation}
Twice $U$ is
$(u_1-1)^2+\sum_{i<k}(u_i-u_{i+1})^2+(u_k-1)^2$.
It is nonnegative, and its zeros make every unit alias equal to one.

\section{Exact coefficient collection replaces a global restriction}
\label{sec:quotient}

The earlier factor-sharing compiler keeps a bipartite factor graph without
four-cycles \cite{graph127}. That condition ensures that every product of
two source-edge monomials determines its two occurrences uniquely. It is
a convenient sufficient condition, but it also restricts products which
never appear in the graph energy.

For instance, take four pairwise disjoint source edges with monomials
$as,at,bs,bt$. Their factor graph contains a four-cycle, yet their actual
energy is
\[
 (as)^2+(at)^2+(bs)^2+(bt)^2.
\]
All coefficients are one. The algebraic equality $(as)(bt)=(at)(bs)$ is
irrelevant because neither cross product occurs. More generally, even
colliding contributions can be harmless when their signed collected
coefficient remains in $\{-1,0,1\}$.

\begin{definition}[Admissible alias identification]
An identification maps old aliases to new aliases only within the same
logical source-variable class, or within the unit class. Flow coordinates
are retained separately. Substitute this map into the graph-edge
polynomials, reconstruct the equality cycles and unit path, and optionally
use Lemma~\ref{lem:pair} for two-alias classes. The identification is
\emph{admissible} if exact expansion shows that every collected coefficient
of the complete resulting energy is $0$, $1$, or $-1$.
\end{definition}

\begin{theorem}[Certified alias quotient]\label{thm:quotient}
Let a normalized integer source system have a certified graph representation
as above. Every admissible alias identification produces a polynomial of
degree at most four and coefficient height at most one, whose integer zeros
are in bijection with the integer zeros of the normalized source. If its
collected list contains a degree-four monomial and constant coefficient one,
its degree and coefficient height are exactly four and one.
\end{theorem}
\begin{proof}
Write the resulting polynomial as
\[
 P'=\frac12\sum_v d_v'^2+\sum_{\text{cycle classes}}C+
          \sum_{\text{pair classes}}F+U.
\]
Every summand is nonnegative over $\R$. Thus $P'=0$ forces every divergence,
every alias equality, and every unit condition. A new alias class therefore
has one well-defined logical source value. Substitute these values into
the source term occurrences. Their logical identities are unchanged by
the identification. Summing the zero divergences over each row block
recovers all normalized source equations.

Conversely, take an integer zero of the normalized source. Assign each
new alias its logical source value and each unit alias the value one.
Every equality and unit energy is zero. The same external row-term values
satisfy all block balances, so Lemma~\ref{lem:tree} supplies the unique
integer internal flows. Consequently $P'=0$.

Every alias and internal flow is uniquely determined by the source zero,
including any retained dummy aliases. Reading source values and performing
this construction are inverse maps. This proves bijection. The degree
bound follows from the degrees of the graph edges and equality energies.
The coefficient assertion is precisely the exact finite admissibility
check. A literal degree-four coefficient and the constant one witness
attainment of both bounds.
\end{proof}

\begin{remark}[Zero sets and polynomial values]
The theorem does not claim the off-solution identity
$P'(L(z))=\sum_r R_r(z)^2$. Canonical rooted-tree elimination leaves one
residual per row block and $-\sum_rR_r$ at the hub root; its graph energy
is instead
\[
 \tfrac12\left(\sum_rR_r^2+(\sum_rR_r)^2\right).
\]
Both have the same zero set. Distinguishing these identities prevents an
incorrect claim of value-preserving substitution.
\end{remark}

\subsection{How the particular certificates were found}

The search merges two aliases of the same logical value and updates all
affected collected monomials exactly. A proposed merge is accepted only if
all new coefficients stay in $\{-1,0,1\}$. Equality cycles are rebuilt
on the distinct surviving aliases. Each safe two-coordinate quartic
replacement is then checked against the full polynomial. All computation
uses integers, with no floating-point comparisons or random evaluations
as proof substitutes.

The order of attempted merges is an optimization choice. A different order
can give a different valid quotient; neither the deterministic order nor
the finite hub sweep proves optimality. The accepted certificate is checked
independently of the search history. This is essential: trust in a greedy
search is unnecessary once its exact finite output has been verified.

The parent graph, the map of source occurrences, the parent-to-child alias
map, all new equality classes, every edge, and the final coefficient list
are retained. The independent checker expands twice the energy as a sum
of squares and compares every coefficient. Its expansion does not use the
producer's adjacent-edge collection routine.

\clearpage
\section{A direct quartic in 306 unrestricted integers}
\label{sec:signed}
\label{sec:signed306}

The main direct construction starts from the literal centered polynomial
of package 131~\cite{cap131}, in 324 unrestricted integer coordinates.  The source is
pinned to ProveIt commit
\srcpin; its equation, centering
certificate, and ordinary arithmetic proof are retained in the archive.
The new polynomial is
\begin{equation}\label{eq:306definition}
 P_{306}(z_0,\ldots,z_{305})=\sum_{j=0}^{205}r_j(z)^2,
 \qquad z\in\mathbb Z^{306}.
\end{equation}
Every $r_j$ is an explicitly listed integer polynomial of degree at most two.  There
are no implicit domain predicates: all 306 coordinates are unrestricted.
The literal residuals and the complete collected coefficient list are
supplied, together with their independent verifier.

\begin{center}
\begin{tabular}{lrrrrr}
\toprule
Equation & Variables & Rows & Degree & Terms & Height\\
\midrule
Pinned parent &324&208&4&5808&128446\\
New equation &306&206&4&6259&128446\\
\bottomrule
\end{tabular}
\end{center}

Thus 18 coordinates and two rows are removed at the cost of 451 collected
terms; both coefficient height and residual height, respectively 128446
and 256, are preserved.  The comparison is within this direct C222
construction.  Much smaller variable counts exist in the repository with
different degrees and coefficient architectures.

\begin{theorem}\label{thm:306-equivalence}
Let $P_{324}$ be the pinned parent polynomial.  Ordinary mathematics
formalizable in ZFC proves
\[
 (\exists z\in\mathbb Z^{306}\;P_{306}(z)=0)
 \quad\Longleftrightarrow\quad
 (\exists w\in\mathbb Z^{324}\;P_{324}(w)=0).
\]
Consequently, composing with the retained parent calibration,
\[
 \mathrm{ZFC}\vdash
 \bigl[\neg\exists z\in\mathbb Z^{306}\;P_{306}(z)=0\bigr]
 \longleftrightarrow\operatorname{Con}(\mathrm{ZFC}).
\]
\end{theorem}

The new idea is a linear orientation term inside the existing bounded
Pell-index congruence.  It permits nine sign graphs to be replaced by one
strict product.  The proof must recover the signs before invoking the
parent's complete arithmetic theorem.  We give that order explicitly.

\subsection{Linear orientation in a bounded residue pool}

For $A\geq2$ let $X_j(A),Y_j(A)$ be the canonical Pell sequences:
\[
 X_0=1,\quad Y_0=0,\qquad
 X_{j+1}=AX_j+(A^2-1)Y_j,\quad Y_{j+1}=X_j+AY_j.
\]
The unchanged parent square frame and complete anchor yield, independently
of the nine encoded-coordinate signs,
\begin{equation}\label{eq:306frame}
 W\geq258,\qquad A>W^W,\qquad A-1>4W^3.
\end{equation}
They also give 15 principal Pell norms, absolute indices $0\leq j_i<W$,
signed proposed indices $|k_i|<W$, and absolute bounds $|n|,|m|<W$ for
every requested base and output.  The fully indexed anchor is at $W$;
its unchanged cap bounds every principal ordinate in absolute value.
Its proof uses the frame, the anchor common-multiplier graph and the cap
graph, all retained here.  No $k_i\geq0$ premise is required.

Pell classification supplies signs $\sigma_i,\epsilon_i\in\{1,-1\}$ with
\[
 x_i=\sigma_iX_{j_i}(A),\qquad y_i=\epsilon_iY_{j_i}(A).
\]
At $j_i=0$ the choice of $\epsilon_i$ is immaterial.  Modulo $A-1$,
the recurrence gives $X_j\equiv1$ and $Y_j\equiv j$.

\begin{lemma}[Linear orientation pool]\label{lem:306pool}
Assume the preceding bounds.  Let $I$ be a set of nine of the fifteen
principal pairs.  The single row
\begin{equation}\label{eq:306pool}
 \sum_{i=1}^{15}(y_i-k_i)^2+\sum_{i\in I}(1-x_i)=(A-1)q_0,
 \qquad q_0\in\mathbb Z,
\end{equation}
forces $\epsilon_i j_i=k_i$ for every $i$ and $\sigma_i=1$ for $i\in I$.
\end{lemma}

\begin{proof}
Reduce \eqref{eq:306pool} modulo $A-1$.  It makes $A-1$ divide
\[
 E=\sum_{i=1}^{15}(\epsilon_i j_i-k_i)^2+
     \sum_{i\in I}(1-\sigma_i).
\]
Every term in the last sum belongs to $\{0,2\}$.  Therefore
\[
 0\leq E\leq60(W-1)^2+18<4W^3<A-1.
\]
For example,
$4W^3-60(W-1)^2-18=4W^2(W-15)+120W-78>0$ for $W\geq258$.
Hence $E=0$, and every nonnegative summand vanishes.
\end{proof}

Only the residue representative $E$ is small and nonnegative.  The actual
left side of \eqref{eq:306pool} need not be either.  In particular the pool
quotient must remain unrestricted.  A squared term $(x_i-1)^2$ was used
in the parent's conservative orientation variant; the linear term
$1-x_i$ works because its residue has the nonnegative alphabet $\{0,2\}$.
The new application combines this cheaper orientation with removal of
the nine encoded-coordinate sign graphs.

\subsection{Signed indices and the early bootstrap}

We need a power lemma that treats zero and negative proposed indices.
Write $Y_{-j}=-Y_j$ and $X_{-j}=X_j$.

\begin{lemma}[Signed-index power request]\label{lem:306request}
Let $2\leq n<W$, $|k|<W$, $|m|<W$, and put
$T=2An-n^2-1$.  Suppose
\[
 y=Y_k(A),\qquad x=\pm X_{|k|}(A),\qquad
 x+(n-A)y=m+Tz
\]
for an integer $z$.  If $x=+X_{|k|}$, then $k\geq0$ and $m=n^k$.
If $k>0$, either choice initially allowed for $x$ is forced to be positive.
At $k=0$ and positive $x$, the output is exactly $m=1$.
\end{lemma}

\begin{proof}
For $j\geq0$ put
$H_j=X_j+(n-A)Y_j$ and $J_j=X_j-(n-A)Y_j$.  Direct multiplication and
the recurrence give
\[
 H_jJ_j=1+TY_j^2,\qquad H_{j+1}=nH_j+TY_j,
\]
so $H_j\equiv n^j\pmod T$.  The frame bounds imply $T>W^W$.  For
$1\leq j<W$ the required differences have absolute value below $W^W$;
in particular
$|mn^j-1|\leq(W-1)^W+1<W^W$.

If $k=-j<0$ and $x=+X_j$, the request gives
$mn^j\equiv1\pmod T$, hence $mn^j=1$ as an integer equality.  This is
impossible since $n^j\geq2$.  If $k=j>0$ and $x=-X_j$, it similarly gives
$mn^j=-1$, also impossible.  The remaining positive-index branch gives
$m\equiv n^k\pmod T$, and the size bound makes it equality.  At $k=0$,
$x=1,y=0$ gives $m=1$ directly by the same small-residue argument.
\end{proof}

The selected nine pairs are precisely the requests at base $q+1$ for
the eight encoded CRT columns and their reciprocal-index column.  The
following bootstrap does not use any of their removed sign graphs.
The retained rows and retained nonnegative quantities $r,s,o$ give
\[
 t=16+h_0^2\geq16,\qquad B=r+s+1\geq1,\qquad q=B+o+1\geq2.
\]
The t-pair has its signed intended index from Lemma~\ref{lem:306pool}.
Its genuine request at base $q$ and positive index $t$ forces its abscissa
positive by Lemma~\ref{lem:306request}.  The same pair therefore gives
\[
 V_0=(F^2+1)^t,
\]
including the possible base-one case before the sign of $F$ is known.
For that case the same recurrence gives $H_t\equiv1\pmod{2(A-1)}$;
the bound $|V_0|<W$ then forces $V_0=1$.
The literal ceiling rows give $FB=V_0+s>0$, so $F\geq1$.  With
$b=F^2+1$, both $b>F$ and $b^{t-1}>F$, whence
\begin{equation}\label{eq:306Bbound}
 FB\geq b^t>Fb,\qquad B>F^2+1.
\end{equation}
The unchanged cofinal row $y=V_0(V_0-1)$ also gives $y>F$; the x-sign
graph remains present.  Neither factorial meaning nor CRT interpretation
has been used at this stage.

Lemma~\ref{lem:306pool} orients the nine selected abscissas, and
Lemma~\ref{lem:306request}, at base $q+1\geq3$, gives
\[
 k_a,k_b,k_c,k_d,k_S,k_{S'},k_u,k_{v_q},k_{\mathrm{index}}\geq0.
\]
Zero remains a genuine possibility at this point.  The reciprocal-index
rows give
\[
 B\,v=F^2+k_{\mathrm{index}}+1
\]
for an unrestricted integer $v$.  If $k_{\mathrm{index}}=0$, its right
side lies strictly between zero and $B$ by \eqref{eq:306Bbound}, an
impossible integer multiple of $B$.  Therefore $k_{\mathrm{index}}\geq1$.

\subsection{Six products and one fused positivity row}

The other eight indices are nonnegative.  Introduce six unrestricted
coordinates $P_2,\ldots,P_7$ with
\begin{align*}
 P_2&=k_a k_b,&P_3&=P_2 k_c,&P_4&=P_3 k_d,\\
 P_5&=P_4 k_S,&P_6&=P_5 k_{S'},&P_7&=P_6 k_u.
\end{align*}
Finally introduce three unrestricted integers $a,b,c$ and impose
\begin{equation}\label{eq:306product}
 P_7 k_{v_q}-1=a^2+b^2+c(c+1).
\end{equation}
The right side is nonnegative at every integer triple.  The product of
the eight nonnegative indices is thus at least one, so every index is
at least one.  This restores all nine old encoded-coordinate signs.
Every displayed row is quadratic; a separate $P_8$ coordinate is avoided.

Conversely every nonnegative integer $n$ is
$a^2+b^2+c(c+1)$.  Indeed, $4n+1$ is not of the forbidden form
$4^u(8v+7)$, so Legendre's three-square theorem represents it by three
squares.  Exactly one is odd; writing the three terms as
$(2a)^2,(2b)^2,(2c+1)^2$ gives the result.  Thus every old zero, whose
eight indices are positive, has witnesses for \eqref{eq:306product}.

The eight-factor idea has a predecessor in the repository's
\texttt{signed-pell-cap-refined-scalar-product} proof.  That proof obtains
nonnegative indices from a different residue-one chart.  Here the original
raw remainder chart is retained; the linear orientation pool supplies the
required signs directly.  The change is this complete composition and its
literal smaller polynomial, not invention of positive-product reasoning.

\subsection{Literal chart and exact witness transformations}

The centered parent has the following requested-index chart.  The third
column lists its nine selected abscissa coordinates; all have zero
centering shift.
\begin{center}
\begin{tabular}{lrr}
\toprule
Index & Expression & Abscissa coordinate\\
\midrule
$k_a$ &$z_{10}-1$&174\\
$k_b$ &$z_{11}$&180\\
$k_c$ &$z_{12}$&184\\
$k_d$ &$z_{13}$&188\\
$k_S$ &$z_{14}$&192\\
$k_{S'}$ &$z_{15}$&200\\
$k_u$ &$z_{16}$&204\\
$k_{v_q}$ &$z_{17}$&212\\
$k_{\mathrm{index}}$ &$z_{18}$&216\\
\bottomrule
\end{tabular}
\end{center}
The old nonnegative quantities are $k_i-1$.  Delete precisely parent rows
184--192, the nine scalar graphs labeled \texttt{domain\_guard.3} through
\texttt{domain\_guard.11}, and their 27 private coordinates 252--278.
They occur nowhere else.  Modify parent row 180, the pooled-index row, by
adding precisely
\[
 9-z_{174}-z_{180}-z_{184}-z_{188}-z_{192}
   -z_{200}-z_{204}-z_{212}-z_{216}.
\]
The private pool quotient is coordinate 242 and occurs only in this row.
The actual Pell center is $A=z_{151}+1$, so its divisor is $A-1=z_{151}$.

For a new zero, the preceding argument proves the nine missing signs.
Let $T_0=\sum_{i\in I}(1-X_i)$.  Each selected abscissa is canonical,
so $T_0/(A-1)$ is an integer.  Set
\[
 q_{\mathrm{old}}=q_{\mathrm{new}}-T_0/(A-1),
\]
and choose fresh Legendre triples for all nine deleted scalar graphs.
Every other retained parent coordinate stays fixed.  The private quotient
change restores the old pooled row and affects no other row.  We now have
a literal zero of the entire parent polynomial, so its complete C222
calibration may be invoked.

For the reverse direction, a parent zero has positive intended indices
and canonical abscissas by its ordinary power theorem.  Set
\[
 q_{\mathrm{new}}=q_{\mathrm{old}}+T_0/(A-1),
\]
forget the 27 old guard witnesses, calculate the six product coordinates,
and supply a triple for \eqref{eq:306product}.  All new rows hold.  These
maps preserve the native front and every retained parent coordinate except
the private pool quotient.  The Legendre witnesses are not unique, so
there is no claim of a bijection between individual full zero tuples.
This proves Theorem~\ref{thm:306-equivalence}.

The parent projection and centering are also part of the supplied
certificate boundary.  The nine deleted guard rows survived projection;
no restoration polynomial depends on their removed auxiliary coordinates
or the private pool quotient.  The independently checked occurrence sets
therefore justify using the parent frame and anchor before restoring those
graphs, without a circular appeal to the entire parent theorem.

\subsection{Finite certification and its scope}

There are $324-27+6+3=306$ coordinates and $208-9+6+1=206$ residuals.
Their sum of squares has a nonzero degree-four coefficient, so the degree
is exactly four.  The independent verifier imports no producer code.  It
checks every retained row, the complete changed pool row, every removed
witness occurrence set, the quotient's privacy, the affine chart, the
six products, the fused positivity row, all collected coefficients and
the exact metrics.  It also rejects deliberate changes to an orientation
term, a factor shift, quotient privacy, and an expanded coefficient.

These finite checks certify literal algebra and artifact bindings.  They
do not establish unbounded number theory by sampling.  Pell classification,
the retained full anchor, Legendre's theorem and the ordinary parent
C222/ZF/ZFC calibration remain the identified mathematical dependencies.
No new end-to-end proof-assistant certificate or encoded object-ZFC proof
is claimed.


\section{Logical calibration and what is inherited}
\label{sec:calibration}

The direct source is the ordinary numerical proof-calculus construction
called C222 in ProveIt. The relevant logical interfaces are
\path{KEndpoint.c222_iff_zf_refutation} and
\path{ConstructibleChoice.zf_refutation_iff_zfc_refutation}.
The arithmetic source relates an integer zero of its explicit residual
system to a finite C222 derivation. The cited interfaces then relate such
a derivation to a finite syntactic ZFC refutation. These are inherited
calibration results, with the precise source snapshot and local retained
proofs recorded in the provenance ledger.

The new steps in this report are the shared orientation/sign construction,
the exact alias-quotient criterion and its realized contractions, and the
two-coordinate equality energy as used in that compiler. They use the
already supplied source arithmetic calibration. A correctness theorem for
a compiler or an exact checksum is not itself a proof of that calibration;
our verification reports keep these roles separate.

All new arguments are finite algebra, integer congruences, elementary
Pell recurrence arguments, finite graph induction, and explicitly inherited
number-theoretic representation theorems. They can be expressed and proved
in ZFC. Composing them with the inherited calibration gives
\[
 \ZFC\vdash
   (\exists z\;P_{\mathrm{dir}}(z)=0)
   \longleftrightarrow
   (\exists p\;\operatorname{Refutation}_{\ZFC}(p)),
\]
and the graph theorem gives the same statement for each unit-coefficient
equation. Classical negation yields \eqref{eq:main}.

There is no newly generated Lean or Rocq proof term for the entire final
coefficient file, and no claim that a newly emitted object-language ZFC
derivation has been checked by a kernel. The release provides an ordinary
proof with exact finite arithmetic certificates. Its mathematical conclusion
does not depend on assuming consistency; interpreting the equations as
having no integer solution does require the corresponding consistency
claim, which we do not prove here.

The direct branch does not simulate K215 and does not instantiate Jones's
universal polynomial. Those other branches remain useful comparisons, not
requirements of this construction. General universal-pair results such as
\cite{bayer,jones} address other arity/degree tradeoffs and do not imply
that the present polynomials are globally minimal.

\section{Exact artifacts and verification}
\label{sec:verification}

The delivery is designed to make each mathematical claim reviewable as a
finite object. The original parent residual system is retained without
alteration. The direct transformation certificate specifies every deleted
guard, changed pooled row, new product equation, and coordinate renaming.
The graph certificates retain the complete normalized source and graph
interfaces. All expanded polynomials are supplied in sparse integer form.

For the direct equation, the first independent checker verifies 929 literal
conditions and rejects four deliberately altered artifacts: a changed
orientation term, a wrong first-factor shift, an escaping private quotient,
and a changed collected coefficient. A separately written reviewer checks
all 199 retained rows, all seven new rows, the nine actual raw power ports,
and all 216 parent restoration polynomials. None of those restorations
depends on the removed private guard coordinates or the changed private
pool quotient. The second reviewer re-expands the full quartic by unordered
monomial pairs, independently of the producer's ordered-pair algorithm.
The coefficient of $z_0^4$ is 85, a literal degree-attainment witness.

For each graph profile the checker first verifies the original normalized
source and graph certificate. It then checks the complete logical alias
map, separate flow coordinates, regenerated equality constraints and unit
path. It reconstructs twice the final polynomial as a sum of squares,
including the three-square identity for each quartic pair guard, and
compares every collected coefficient. It also parses every term of the
readable \path{polynomial.txt} back to the normative coefficient list.

All three graph profiles have been reconstructed from their source rows,
with byte-identical parent polynomials, parent certificates, final
polynomials, final certificates, and metric files: eighteen mathematical
JSON artifacts in total. The optional reconstruction uses the builders only
after the independent verification. Normal and optimization-mode Python
runs retain all necessary checks; the verifiers do not depend on
\texttt{assert} statements.
The integrated reconstruction also recreates all six direct-equation files
in a temporary directory and compares their exact bytes.

The graph evaluator also supplies a canonical witness lift. For an arbitrary
integer source tuple $z$, its construction satisfies the exact identity
\[
 2P_{\mathrm{unit}}(L(z))
   =\sum_rR_r(z)^2+\left(\sum_rR_r(z)\right)^2.
\]
On source zeros this becomes a constructive lift to graph zeros. The
identity follows from the block-tree construction and the monic coefficient
normalization; finite evaluations are supplementary diagnostics.

The integrated release check additionally compares the 306-variable source
copied into the graph component with the direct equation byte for byte,
parses every printed LaTeX residual back into its exact coefficient map,
and checks the equation identities and integrity inventory. Its receipt is
\path{verification/RELEASE_CHECK.json}; the optimization-mode companion is
\path{verification/RELEASE_CHECK_OPTIMIZED.json}. Neither check is a search
for a zero or a finite test of consistency.
The retained build receipts precede final inventory sealing; running the
released checker also verifies the completed \path{MANIFEST.sha256}.


The distinction between a mathematical proof and an integrity check matters.
A SHA-256 inventory detects changed bytes; it does not establish an
algebraic identity. Exact coefficient reconstruction proves the stated
finite identity; it does not prove the unbounded Pell or logical
calibration by testing examples. The release includes these layers
separately and records which programs check which assertions.

\subsection{Offline use}

Python 3.10 or later with its standard library suffices for the selected
arithmetic-artifact checks. From the release directory run
\begin{lstlisting}
python3 -B verify.py
python3 -B -O verify.py
python3 -B verify.py --rebuild
python3 -B evaluate.py --help
\end{lstlisting}
The verification entry point preserves the delivered equations. Evaluation
at a supplied integer tuple is exact, but is not an algorithm for deciding
whether any integer zero exists. The ordinary proof and inherited
calibration give the semantic interpretation of the polynomial; the
evaluator simply evaluates its literal coefficients.

The paper is built with
\begin{lstlisting}
latexmk -pdf -interaction=nonstopmode -halt-on-error report.tex
\end{lstlisting}
The LaTeX inputs, generated tables, and readable equation transcriptions are
included. PDF byte identity can depend on TeX versions and metadata; the
mathematical coefficient lists and their canonical reconstruction are the
normative finite artifacts.

\section{Further research questions}
\label{sec:future}

\paragraph{Joint choice of graph and aliases.}
The hub layout affects which edge products occur, and hence which alias
identifications are safe. The current search chooses a graph first and
contracts aliases afterwards. A joint optimizer could trade a small number
of tree-flow coordinates for many fewer factor aliases, or reduce support
without increasing arity. A graph modification would need a new block-tree
certificate, not only a coefficient count.

\paragraph{Larger certified contractions.}
A greedy sequence of individually admissible pair merges can miss a
simultaneous merge whose intermediate polynomials have height two but whose
final signed collisions cancel back to height one. Theorem~\ref{thm:quotient}
applies directly to the final identification. Searching blocks of aliases
with an independently checked final certificate is therefore legitimate.
There is no need to insist that every hypothetical intermediate quotient
itself be released as a height-one polynomial.

\paragraph{Other equality energies.}
The quartic identity in Lemma~\ref{lem:pair} is useful but collides with
fourth powers already present in some graph energies. Can a small library
of nonnegative height-one equality polynomials, used with exact overlap
tests, eliminate more dummies? Any alternative must preserve its full zero
set over unrestricted integers, including zero and negative values. An
individual-degree-two companion would be especially useful, since the
present quartic pair guard permits fourth powers.

\paragraph{A general finite-residue orientation compiler.}
The Pell trace has the tiny residue alphabet $\{-1,1\}$ modulo $A-1$.
The new linear defect $1-x$ represents the nonnegative values $0$ and $2$
on that alphabet. Other arithmetic modules may have similarly small
residue alphabets, permitting their branch choices to be certified in one
pooled congruence. The required error bound and dependency order must be
proved before any sign information is used.

\paragraph{Formal certificates for the entire chain.}
The graph and alias theorems are natural targets for Lean or Rocq. A small
verified sparse-polynomial checker could then connect their abstract
statements to the emitted coefficient files. The direct orientation lemma
needs a formal signed-Pell classification and the exact source frame bound.
Closing these interfaces would strengthen the verification status without
changing the literal equations.

\paragraph{Comparisons under a stated objective.}
It remains reasonable to minimize arity subject to degree four and height
one, or minimize support subject to a fixed arity cap. These are concrete
optimization problems. A single unqualified phrase such as ``the simplest
Diophantine equation'' does not select between such incomparable outcomes.
The present release makes progress under explicit criteria and leaves the
global question open.

\appendix
\section{Literal direct residuals}
\label{app:residuals}
Every coordinate in this appendix ranges independently over $\Z$. The
single polynomial is the sum of the squares of \emph{all} the displayed
residuals, not the square of their sum. Coordinate names impose no extra
domain restrictions. This transcription is generated from the same literal
JSON checked by the independent expansion verifier.
\par\begingroup\small
\setlength{\abovedisplayskip}{4pt}\setlength{\belowdisplayskip}{5pt}
\par\noindent\begin{minipage}{\linewidth}
\noindent{\scriptsize\textsf{0: multiplication}}
\begin{align*}
r_{0}={}&-z_{3}^{2}+z_{4}
\end{align*}
\end{minipage}\par
\par\noindent\begin{minipage}{\linewidth}
\noindent{\scriptsize\textsf{1: multiplication}}
\begin{align*}
r_{1}={}&-1-z_{0}z_{2}+z_{9}
\end{align*}
\end{minipage}\par
\par\noindent\begin{minipage}{\linewidth}
\noindent{\scriptsize\textsf{2: multiplication}}
\begin{align*}
r_{2}={}&-1+z_{1}+z_{2}-z_{2}z_{7}-z_{3}-z_{7}
\end{align*}
\end{minipage}\par
\par\noindent\begin{minipage}{\linewidth}
\noindent{\scriptsize\textsf{3: multiplication}}
\begin{align*}
r_{3}={}&-2+z_{1}-z_{8}z_{9}
\end{align*}
\end{minipage}\par
\par\noindent\begin{minipage}{\linewidth}
\noindent{\scriptsize\textsf{4: multiplication}}
\begin{align*}
r_{4}={}&1+z_{4}-z_{5}z_{20}+z_{18}
\end{align*}
\end{minipage}\par
\par\noindent\begin{minipage}{\linewidth}
\noindent{\scriptsize\textsf{5: multiplication}}
\begin{align*}
r_{5}={}&-1+2z_{10}-z_{10}^{2}-2z_{10}z_{11}+2z_{11}-z_{11}^{2}+z_{21}
\end{align*}
\end{minipage}\par
\par\noindent\begin{minipage}{\linewidth}
\noindent{\scriptsize\textsf{6: multiplication}}
\begin{align*}
r_{6}={}&-5-30z_{10}-36z_{10}^{2}-96z_{10}z_{11}-24z_{10}z_{21}-40z_{11}-64z_{11}^{2}-32z_{11}z_{21}-10z_{21}-4z_{21}^{2}\\[-2pt]
&+z_{22}
\end{align*}
\end{minipage}\par
\par\noindent\begin{minipage}{\linewidth}
\noindent{\scriptsize\textsf{7: multiplication}}
\begin{align*}
r_{7}={}&-2+2z_{10}-z_{10}^{2}-4z_{10}z_{22}+4z_{22}-4z_{22}^{2}+z_{23}
\end{align*}
\end{minipage}\par
\par\noindent\begin{minipage}{\linewidth}
\noindent{\scriptsize\textsf{8: multiplication}}
\begin{align*}
r_{8}={}&1+6z_{10}-36z_{10}^{2}-192z_{10}z_{22}-24z_{10}z_{23}+16z_{22}-256z_{22}^{2}-64z_{22}z_{23}+2z_{23}-4z_{23}^{2}+z_{24}
\end{align*}
\end{minipage}\par
\par\noindent\begin{minipage}{\linewidth}
\noindent{\scriptsize\textsf{9: multiplication}}
\begin{align*}
r_{9}={}&-2-z_{11}^{2}-2z_{11}z_{12}-z_{12}^{2}+z_{25}
\end{align*}
\end{minipage}\par
\par\noindent\begin{minipage}{\linewidth}
\noindent{\scriptsize\textsf{10: multiplication}}
\begin{align*}
r_{10}={}&-9+7z_{10}-z_{10}^{2}-4z_{10}z_{11}-4z_{10}z_{25}+12z_{11}-4z_{11}^{2}-8z_{11}z_{25}+12z_{25}-4z_{25}^{2}+z_{26}
\end{align*}
\end{minipage}\par
\par\noindent\begin{minipage}{\linewidth}
\noindent{\scriptsize\textsf{11: multiplication}}
\begin{align*}
r_{11}={}&-1+2z_{10}-z_{10}^{2}-2z_{10}z_{12}+2z_{12}-z_{12}^{2}+z_{27}
\end{align*}
\end{minipage}\par
\par\noindent\begin{minipage}{\linewidth}
\noindent{\scriptsize\textsf{12: multiplication}}
\begin{align*}
r_{12}={}&-6z_{10}-144z_{10}^{2}-192z_{10}z_{11}-192z_{10}z_{12}-48z_{10}z_{21}-48z_{10}z_{27}-8z_{11}-64z_{11}^{2}\\[-2pt]
&-128z_{11}z_{12}-32z_{11}z_{21}-32z_{11}z_{27}-64z_{12}^{2}-32z_{12}z_{21}-32z_{12}z_{27}-2z_{21}-4z_{21}^{2}\\[-2pt]
&-8z_{21}z_{27}-4z_{27}^{2}+z_{28}
\end{align*}
\end{minipage}\par
\par\noindent\begin{minipage}{\linewidth}
\noindent{\scriptsize\textsf{13: multiplication}}
\begin{align*}
r_{13}={}&-4z_{26}^{2}-8z_{26}z_{28}-4z_{28}^{2}+z_{29}
\end{align*}
\end{minipage}\par
\par\noindent\begin{minipage}{\linewidth}
\noindent{\scriptsize\textsf{14: multiplication}}
\begin{align*}
r_{14}={}&-6+16z_{12}-16z_{12}^{2}-16z_{12}z_{25}-16z_{12}z_{27}+8z_{25}-4z_{25}^{2}-8z_{25}z_{27}+8z_{27}-4z_{27}^{2}+z_{30}
\end{align*}
\end{minipage}\par
\par\noindent\begin{minipage}{\linewidth}
\noindent{\scriptsize\textsf{15: multiplication}}
\begin{align*}
r_{15}={}&1+6z_{10}-36z_{10}^{2}-96z_{10}z_{11}-48z_{10}z_{12}-24z_{10}z_{21}-48z_{10}z_{27}-24z_{10}z_{30}+8z_{11}-64z_{11}^{2}\\[-2pt]
&-64z_{11}z_{12}-32z_{11}z_{21}-64z_{11}z_{27}-32z_{11}z_{30}+8z_{12}-16z_{12}^{2}-16z_{12}z_{21}-32z_{12}z_{27}\\[-2pt]
&-16z_{12}z_{30}+2z_{21}-4z_{21}^{2}-16z_{21}z_{27}-8z_{21}z_{30}+8z_{27}-16z_{27}^{2}-16z_{27}z_{30}+4z_{30}\\[-2pt]
&-4z_{30}^{2}+z_{31}
\end{align*}
\end{minipage}\par
\par\noindent\begin{minipage}{\linewidth}
\noindent{\scriptsize\textsf{16: multiplication}}
\begin{align*}
r_{16}={}&-10-5z_{10}-z_{10}^{2}-12z_{10}z_{11}-16z_{10}z_{12}-4z_{10}z_{25}-36z_{11}-36z_{11}^{2}-96z_{11}z_{12}-24z_{11}z_{25}\\[-2pt]
&-48z_{12}-64z_{12}^{2}-32z_{12}z_{25}-12z_{25}-4z_{25}^{2}+z_{32}
\end{align*}
\end{minipage}\par
\par\noindent\begin{minipage}{\linewidth}
\noindent{\scriptsize\textsf{17: multiplication}}
\begin{align*}
r_{17}={}&-2-24z_{10}-36z_{10}^{2}-12z_{10}z_{11}-96z_{10}z_{12}-24z_{10}z_{27}-3z_{11}-z_{11}^{2}-16z_{11}z_{12}-4z_{11}z_{27}\\[-2pt]
&-32z_{12}-64z_{12}^{2}-32z_{12}z_{27}-8z_{27}-4z_{27}^{2}+z_{33}
\end{align*}
\end{minipage}\par
\par\noindent\begin{minipage}{\linewidth}
\noindent{\scriptsize\textsf{18: multiplication}}
\begin{align*}
r_{18}={}&2-4z_{32}^{2}-8z_{32}z_{33}-4z_{33}^{2}+z_{34}
\end{align*}
\end{minipage}\par
\par\noindent\begin{minipage}{\linewidth}
\noindent{\scriptsize\textsf{19: multiplication}}
\begin{align*}
r_{19}={}&-12+16z_{10}-4z_{10}^{2}-8z_{10}z_{12}-8z_{10}z_{21}-8z_{10}z_{25}+16z_{12}-4z_{12}^{2}-8z_{12}z_{21}-8z_{12}z_{25}\\[-2pt]
&+16z_{21}-4z_{21}^{2}-8z_{21}z_{25}+16z_{25}-4z_{25}^{2}+z_{35}
\end{align*}
\end{minipage}\par
\par\noindent\begin{minipage}{\linewidth}
\noindent{\scriptsize\textsf{20: multiplication}}
\begin{align*}
r_{20}={}&-3-24z_{10}-64z_{10}^{2}-160z_{10}z_{12}-64z_{10}z_{25}-32z_{10}z_{27}-32z_{10}z_{35}-26z_{12}-100z_{12}^{2}\\[-2pt]
&-80z_{12}z_{25}-40z_{12}z_{27}-40z_{12}z_{35}-8z_{25}-16z_{25}^{2}-16z_{25}z_{27}-16z_{25}z_{35}-6z_{27}-4z_{27}^{2}\\[-2pt]
&-8z_{27}z_{35}-4z_{35}-4z_{35}^{2}+z_{36}
\end{align*}
\end{minipage}\par
\par\noindent\begin{minipage}{\linewidth}
\noindent{\scriptsize\textsf{21: multiplication}}
\begin{align*}
r_{21}={}&-1-18z_{10}-36z_{10}^{2}-96z_{10}z_{11}-24z_{10}z_{21}-24z_{11}-64z_{11}^{2}-32z_{11}z_{21}-6z_{21}-4z_{21}^{2}+z_{37}
\end{align*}
\end{minipage}\par
\par\noindent\begin{minipage}{\linewidth}
\noindent{\scriptsize\textsf{22: multiplication}}
\begin{align*}
r_{22}={}&-2+2z_{10}-z_{10}^{2}-4z_{10}z_{37}+4z_{37}-4z_{37}^{2}+z_{38}
\end{align*}
\end{minipage}\par
\par\noindent\begin{minipage}{\linewidth}
\noindent{\scriptsize\textsf{23: multiplication}}
\begin{align*}
r_{23}={}&1+6z_{10}-36z_{10}^{2}-192z_{10}z_{37}-24z_{10}z_{38}+16z_{37}-256z_{37}^{2}-64z_{37}z_{38}+2z_{38}-4z_{38}^{2}+z_{39}
\end{align*}
\end{minipage}\par
\par\noindent\begin{minipage}{\linewidth}
\noindent{\scriptsize\textsf{24: multiplication}}
\begin{align*}
r_{24}={}&-z_{11}^{2}-4z_{11}z_{39}-4z_{39}^{2}+z_{40}
\end{align*}
\end{minipage}\par
\par\noindent\begin{minipage}{\linewidth}
\noindent{\scriptsize\textsf{25: multiplication}}
\begin{align*}
r_{25}={}&-6-30z_{10}-64z_{10}^{2}-128z_{10}z_{11}-64z_{10}z_{21}-32z_{11}-64z_{11}^{2}-64z_{11}z_{21}-14z_{21}-16z_{21}^{2}\\[-2pt]
&+z_{41}
\end{align*}
\end{minipage}\par
\par\noindent\begin{minipage}{\linewidth}
\noindent{\scriptsize\textsf{26: multiplication}}
\begin{align*}
r_{26}={}&-4+8z_{11}-64z_{11}^{2}-96z_{11}z_{12}-32z_{11}z_{25}-32z_{11}z_{41}+6z_{12}-36z_{12}^{2}-24z_{12}z_{25}\\[-2pt]
&-24z_{12}z_{41}+2z_{25}-4z_{25}^{2}-8z_{25}z_{41}+4z_{41}-4z_{41}^{2}+z_{42}
\end{align*}
\end{minipage}\par
\par\noindent\begin{minipage}{\linewidth}
\noindent{\scriptsize\textsf{27: multiplication}}
\begin{align*}
r_{27}={}&8z_{10}-64z_{10}^{2}-96z_{10}z_{12}-32z_{10}z_{27}-32z_{10}z_{42}+6z_{12}-36z_{12}^{2}-24z_{12}z_{27}-24z_{12}z_{42}\\[-2pt]
&+2z_{27}-4z_{27}^{2}-8z_{27}z_{42}+4z_{42}-4z_{42}^{2}+z_{43}
\end{align*}
\end{minipage}\par
\par\noindent\begin{minipage}{\linewidth}
\noindent{\scriptsize\textsf{28: multiplication}}
\begin{align*}
r_{28}={}&-z_{12}^{2}-2z_{12}z_{13}-z_{13}^{2}+z_{44}
\end{align*}
\end{minipage}\par
\par\noindent\begin{minipage}{\linewidth}
\noindent{\scriptsize\textsf{29: multiplication}}
\begin{align*}
r_{29}={}&2z_{10}-z_{10}^{2}-4z_{10}z_{12}-4z_{10}z_{44}+4z_{12}-4z_{12}^{2}-8z_{12}z_{44}+4z_{44}-4z_{44}^{2}+z_{45}
\end{align*}
\end{minipage}\par
\par\noindent\begin{minipage}{\linewidth}
\noindent{\scriptsize\textsf{30: multiplication}}
\begin{align*}
r_{30}={}&-35+24z_{10}-4z_{10}^{2}-8z_{10}z_{11}-8z_{10}z_{25}-8z_{10}z_{27}+24z_{11}-4z_{11}^{2}-8z_{11}z_{25}-8z_{11}z_{27}\\[-2pt]
&+24z_{25}-4z_{25}^{2}-8z_{25}z_{27}+24z_{27}-4z_{27}^{2}+z_{46}
\end{align*}
\end{minipage}\par
\par\noindent\begin{minipage}{\linewidth}
\noindent{\scriptsize\textsf{31: multiplication}}
\begin{align*}
r_{31}={}&-16z_{10}^{2}-32z_{10}z_{12}-32z_{10}z_{27}-32z_{10}z_{44}-16z_{10}z_{45}-16z_{10}z_{46}-4z_{12}-16z_{12}^{2}\\[-2pt]
&-32z_{12}z_{27}-32z_{12}z_{44}-16z_{12}z_{45}-16z_{12}z_{46}-16z_{27}^{2}-32z_{27}z_{44}-16z_{27}z_{45}-16z_{27}z_{46}\\[-2pt]
&-4z_{44}-16z_{44}^{2}-16z_{44}z_{45}-16z_{44}z_{46}-2z_{45}-4z_{45}^{2}-8z_{45}z_{46}-4z_{46}^{2}+z_{47}
\end{align*}
\end{minipage}\par
\par\noindent\begin{minipage}{\linewidth}
\noindent{\scriptsize\textsf{32: multiplication}}
\begin{align*}
r_{32}={}&z_{12}-25z_{12}^{2}-40z_{12}z_{27}-20z_{12}z_{30}-16z_{27}^{2}-16z_{27}z_{30}-4z_{30}^{2}+z_{48}
\end{align*}
\end{minipage}\par
\par\noindent\begin{minipage}{\linewidth}
\noindent{\scriptsize\textsf{33: multiplication}}
\begin{align*}
r_{33}={}&-1+8z_{11}-4z_{11}^{2}-8z_{11}z_{25}+8z_{25}-4z_{25}^{2}+z_{49}
\end{align*}
\end{minipage}\par
\par\noindent\begin{minipage}{\linewidth}
\noindent{\scriptsize\textsf{34: multiplication}}
\begin{align*}
r_{34}={}&-1-16z_{11}^{2}-8z_{11}z_{12}-32z_{11}z_{25}-16z_{11}z_{49}-z_{12}^{2}-8z_{12}z_{25}-4z_{12}z_{49}-16z_{25}^{2}\\[-2pt]
&-16z_{25}z_{49}-4z_{49}^{2}+z_{50}
\end{align*}
\end{minipage}\par
\par\noindent\begin{minipage}{\linewidth}
\noindent{\scriptsize\textsf{35: multiplication}}
\begin{align*}
r_{35}={}&1+z_{11}-z_{11}^{2}-4z_{11}z_{12}-4z_{11}z_{50}-4z_{12}^{2}-8z_{12}z_{50}-4z_{50}^{2}+z_{51}
\end{align*}
\end{minipage}\par
\par\noindent\begin{minipage}{\linewidth}
\noindent{\scriptsize\textsf{36: multiplication}}
\begin{align*}
r_{36}={}&-1-4z_{51}-4z_{51}^{2}+z_{52}
\end{align*}
\end{minipage}\par
\par\noindent\begin{minipage}{\linewidth}
\noindent{\scriptsize\textsf{37: multiplication}}
\begin{align*}
r_{37}={}&-1-z_{12}^{2}-8z_{12}z_{51}-4z_{12}z_{52}-16z_{51}^{2}-16z_{51}z_{52}-4z_{52}^{2}+z_{53}
\end{align*}
\end{minipage}\par
\par\noindent\begin{minipage}{\linewidth}
\noindent{\scriptsize\textsf{38: multiplication}}
\begin{align*}
r_{38}={}&2+z_{11}-z_{11}^{2}-4z_{11}z_{12}-4z_{11}z_{53}-4z_{12}^{2}-8z_{12}z_{53}-4z_{53}^{2}+z_{54}
\end{align*}
\end{minipage}\par
\par\noindent\begin{minipage}{\linewidth}
\noindent{\scriptsize\textsf{39: multiplication}}
\begin{align*}
r_{39}={}&-4z_{48}^{2}-8z_{48}z_{54}-4z_{54}^{2}+z_{55}
\end{align*}
\end{minipage}\par
\par\noindent\begin{minipage}{\linewidth}
\noindent{\scriptsize\textsf{40: multiplication}}
\begin{align*}
r_{40}={}&-9+7z_{10}-z_{10}^{2}-8z_{10}z_{48}-4z_{10}z_{55}+24z_{48}-16z_{48}^{2}-16z_{48}z_{55}+12z_{55}-4z_{55}^{2}+z_{56}
\end{align*}
\end{minipage}\par
\par\noindent\begin{minipage}{\linewidth}
\noindent{\scriptsize\textsf{41: multiplication}}
\begin{align*}
r_{41}={}&-1-4z_{56}-4z_{56}^{2}+z_{57}
\end{align*}
\end{minipage}\par
\par\noindent\begin{minipage}{\linewidth}
\noindent{\scriptsize\textsf{42: multiplication}}
\begin{align*}
r_{42}={}&-z_{12}^{2}-8z_{12}z_{56}-4z_{12}z_{57}-16z_{56}^{2}-16z_{56}z_{57}-4z_{57}^{2}+z_{58}
\end{align*}
\end{minipage}\par
\par\noindent\begin{minipage}{\linewidth}
\noindent{\scriptsize\textsf{43: multiplication}}
\begin{align*}
r_{43}={}&-z_{10}-z_{10}^{2}-4z_{10}z_{12}-4z_{10}z_{58}-4z_{12}-4z_{12}^{2}-8z_{12}z_{58}-4z_{58}-4z_{58}^{2}+z_{59}
\end{align*}
\end{minipage}\par
\par\noindent\begin{minipage}{\linewidth}
\noindent{\scriptsize\textsf{44: multiplication}}
\begin{align*}
r_{44}={}&-10-12z_{59}-4z_{59}^{2}+z_{60}
\end{align*}
\end{minipage}\par
\par\noindent\begin{minipage}{\linewidth}
\noindent{\scriptsize\textsf{45: multiplication}}
\begin{align*}
r_{45}={}&1+4z_{11}-4z_{11}^{2}-8z_{11}z_{25}-16z_{11}z_{59}-8z_{11}z_{60}+4z_{25}-4z_{25}^{2}-16z_{25}z_{59}-8z_{25}z_{60}\\[-2pt]
&+8z_{59}-16z_{59}^{2}-16z_{59}z_{60}+4z_{60}-4z_{60}^{2}+z_{61}
\end{align*}
\end{minipage}\par
\par\noindent\begin{minipage}{\linewidth}
\noindent{\scriptsize\textsf{46: multiplication}}
\begin{align*}
r_{46}={}&z_{11}-25z_{11}^{2}-40z_{11}z_{25}-20z_{11}z_{61}-16z_{25}^{2}-16z_{25}z_{61}-4z_{61}^{2}+z_{62}
\end{align*}
\end{minipage}\par
\par\noindent\begin{minipage}{\linewidth}
\noindent{\scriptsize\textsf{47: multiplication}}
\begin{align*}
r_{47}={}&-10-12z_{54}-4z_{54}^{2}+z_{63}
\end{align*}
\end{minipage}\par
\par\noindent\begin{minipage}{\linewidth}
\noindent{\scriptsize\textsf{48: multiplication}}
\begin{align*}
r_{48}={}&1-16z_{54}^{2}-16z_{54}z_{62}-16z_{54}z_{63}-4z_{62}^{2}-8z_{62}z_{63}-4z_{63}^{2}+z_{64}
\end{align*}
\end{minipage}\par
\par\noindent\begin{minipage}{\linewidth}
\noindent{\scriptsize\textsf{49: multiplication}}
\begin{align*}
r_{49}={}&z_{12}-z_{12}^{2}-8z_{12}z_{62}-4z_{12}z_{64}-16z_{62}^{2}-16z_{62}z_{64}-4z_{64}^{2}+z_{65}
\end{align*}
\end{minipage}\par
\par\noindent\begin{minipage}{\linewidth}
\noindent{\scriptsize\textsf{50: multiplication}}
\begin{align*}
r_{50}={}&-4z_{10}-4z_{10}^{2}-4z_{10}z_{12}-8z_{10}z_{21}-z_{12}-z_{12}^{2}-4z_{12}z_{21}-4z_{21}-4z_{21}^{2}+z_{66}
\end{align*}
\end{minipage}\par
\par\noindent\begin{minipage}{\linewidth}
\noindent{\scriptsize\textsf{51: multiplication}}
\begin{align*}
r_{51}={}&1+4z_{10}-4z_{10}^{2}-8z_{10}z_{21}-8z_{10}z_{66}+4z_{21}-4z_{21}^{2}-8z_{21}z_{66}+4z_{66}-4z_{66}^{2}+z_{67}
\end{align*}
\end{minipage}\par
\par\noindent\begin{minipage}{\linewidth}
\noindent{\scriptsize\textsf{52: multiplication}}
\begin{align*}
r_{52}={}&-2-4z_{10}-16z_{10}^{2}-32z_{10}z_{21}-16z_{10}z_{65}-16z_{10}z_{67}-4z_{21}-16z_{21}^{2}-16z_{21}z_{65}\\[-2pt]
&-16z_{21}z_{67}-4z_{65}^{2}-8z_{65}z_{67}-2z_{67}-4z_{67}^{2}+z_{68}
\end{align*}
\end{minipage}\par
\par\noindent\begin{minipage}{\linewidth}
\noindent{\scriptsize\textsf{53: multiplication}}
\begin{align*}
r_{53}={}&-2-4z_{68}-4z_{68}^{2}+z_{69}
\end{align*}
\end{minipage}\par
\par\noindent\begin{minipage}{\linewidth}
\noindent{\scriptsize\textsf{54: multiplication}}
\begin{align*}
r_{54}={}&-z_{3}z_{15}+z_{70}
\end{align*}
\end{minipage}\par
\par\noindent\begin{minipage}{\linewidth}
\noindent{\scriptsize\textsf{55: multiplication}}
\begin{align*}
r_{55}={}&z_{14}-z_{15}-2z_{24}z_{70}-z_{71}
\end{align*}
\end{minipage}\par
\par\noindent\begin{minipage}{\linewidth}
\noindent{\scriptsize\textsf{56: multiplication}}
\begin{align*}
r_{56}={}&z_{14}-z_{15}-4z_{26}z_{70}-2z_{29}z_{70}-z_{72}
\end{align*}
\end{minipage}\par
\par\noindent\begin{minipage}{\linewidth}
\noindent{\scriptsize\textsf{57: multiplication}}
\begin{align*}
r_{57}={}&z_{14}-z_{15}-2z_{31}z_{70}-6z_{70}-z_{73}
\end{align*}
\end{minipage}\par
\par\noindent\begin{minipage}{\linewidth}
\noindent{\scriptsize\textsf{58: multiplication}}
\begin{align*}
r_{58}={}&z_{14}-z_{15}-4z_{32}z_{70}-2z_{34}z_{70}-z_{74}
\end{align*}
\end{minipage}\par
\par\noindent\begin{minipage}{\linewidth}
\noindent{\scriptsize\textsf{59: multiplication}}
\begin{align*}
r_{59}={}&z_{14}-z_{15}-2z_{36}z_{70}+6z_{70}-z_{75}
\end{align*}
\end{minipage}\par
\par\noindent\begin{minipage}{\linewidth}
\noindent{\scriptsize\textsf{60: multiplication}}
\begin{align*}
r_{60}={}&z_{14}-z_{15}-4z_{39}z_{70}-2z_{40}z_{70}-z_{76}
\end{align*}
\end{minipage}\par
\par\noindent\begin{minipage}{\linewidth}
\noindent{\scriptsize\textsf{61: multiplication}}
\begin{align*}
r_{61}={}&z_{14}-z_{15}-2z_{43}z_{70}-z_{77}
\end{align*}
\end{minipage}\par
\par\noindent\begin{minipage}{\linewidth}
\noindent{\scriptsize\textsf{62: multiplication}}
\begin{align*}
r_{62}={}&z_{14}-z_{15}-2z_{47}z_{70}-4z_{70}-z_{78}
\end{align*}
\end{minipage}\par
\par\noindent\begin{minipage}{\linewidth}
\noindent{\scriptsize\textsf{63: nonnegative\_square\_difference}}
\begin{align*}
r_{63}={}&-z_{10}z_{11}+z_{10}z_{12}+z_{11}+z_{11}z_{12}-z_{12}-z_{12}^{2}+z_{79}
\end{align*}
\end{minipage}\par
\par\noindent\begin{minipage}{\linewidth}
\noindent{\scriptsize\textsf{64: multiplication}}
\begin{align*}
r_{64}={}&1+2z_{13}+z_{13}^{2}-z_{79}^{2}+z_{80}
\end{align*}
\end{minipage}\par
\par\noindent\begin{minipage}{\linewidth}
\noindent{\scriptsize\textsf{65: nonnegative\_square\_difference}}
\begin{align*}
r_{65}={}&-z_{80}^{2}+z_{81}
\end{align*}
\end{minipage}\par
\par\noindent\begin{minipage}{\linewidth}
\noindent{\scriptsize\textsf{66: multiplication}}
\begin{align*}
r_{66}={}&-4z_{68}z_{70}-2z_{69}z_{70}+2z_{70}+z_{82}
\end{align*}
\end{minipage}\par
\par\noindent\begin{minipage}{\linewidth}
\noindent{\scriptsize\textsf{67: nonnegative\_square\_difference}}
\begin{align*}
r_{67}={}&-z_{14}^{2}+2z_{14}z_{15}+2z_{14}z_{82}-z_{15}^{2}-2z_{15}z_{82}-z_{82}^{2}+z_{83}
\end{align*}
\end{minipage}\par
\par\noindent\begin{minipage}{\linewidth}
\noindent{\scriptsize\textsf{68: multiplication}}
\begin{align*}
r_{68}={}&-1-z_{3}z_{11}+z_{84}
\end{align*}
\end{minipage}\par
\par\noindent\begin{minipage}{\linewidth}
\noindent{\scriptsize\textsf{69: multiplication}}
\begin{align*}
r_{69}={}&-z_{15}z_{84}+z_{86}z_{87}+z_{88}
\end{align*}
\end{minipage}\par
\par\noindent\begin{minipage}{\linewidth}
\noindent{\scriptsize\textsf{70: multiplication}}
\begin{align*}
r_{70}={}&-1+z_{3}-z_{3}z_{10}+z_{85}
\end{align*}
\end{minipage}\par
\par\noindent\begin{minipage}{\linewidth}
\noindent{\scriptsize\textsf{71: multiplication}}
\begin{align*}
r_{71}={}&-1+2z_{3}-2z_{3}z_{10}-2z_{3}z_{21}+z_{86}
\end{align*}
\end{minipage}\par
\par\noindent\begin{minipage}{\linewidth}
\noindent{\scriptsize\textsf{72: multiplication}}
\begin{align*}
r_{72}={}&-z_{14}z_{85}+z_{87}
\end{align*}
\end{minipage}\par
\par\noindent\begin{minipage}{\linewidth}
\noindent{\scriptsize\textsf{73: multiplication}}
\begin{align*}
r_{73}={}&-1-6z_{3}z_{10}-8z_{3}z_{11}-2z_{3}z_{21}+z_{89}
\end{align*}
\end{minipage}\par
\par\noindent\begin{minipage}{\linewidth}
\noindent{\scriptsize\textsf{74: multiplication}}
\begin{align*}
r_{74}={}&z_{14}z_{84}-z_{15}z_{89}+z_{90}
\end{align*}
\end{minipage}\par
\par\noindent\begin{minipage}{\linewidth}
\noindent{\scriptsize\textsf{75: multiplication}}
\begin{align*}
r_{75}={}&-z_{14}z_{71}+z_{15}z_{71}+z_{91}
\end{align*}
\end{minipage}\par
\par\noindent\begin{minipage}{\linewidth}
\noindent{\scriptsize\textsf{76: multiplication}}
\begin{align*}
r_{76}={}&-z_{72}z_{91}+z_{92}
\end{align*}
\end{minipage}\par
\par\noindent\begin{minipage}{\linewidth}
\noindent{\scriptsize\textsf{77: multiplication}}
\begin{align*}
r_{77}={}&-z_{73}z_{92}+z_{93}
\end{align*}
\end{minipage}\par
\par\noindent\begin{minipage}{\linewidth}
\noindent{\scriptsize\textsf{78: multiplication}}
\begin{align*}
r_{78}={}&-z_{74}z_{93}+z_{94}
\end{align*}
\end{minipage}\par
\par\noindent\begin{minipage}{\linewidth}
\noindent{\scriptsize\textsf{79: multiplication}}
\begin{align*}
r_{79}={}&-z_{75}z_{94}+z_{95}
\end{align*}
\end{minipage}\par
\par\noindent\begin{minipage}{\linewidth}
\noindent{\scriptsize\textsf{80: multiplication}}
\begin{align*}
r_{80}={}&-z_{76}z_{95}+z_{96}
\end{align*}
\end{minipage}\par
\par\noindent\begin{minipage}{\linewidth}
\noindent{\scriptsize\textsf{81: multiplication}}
\begin{align*}
r_{81}={}&-z_{77}z_{96}+z_{97}
\end{align*}
\end{minipage}\par
\par\noindent\begin{minipage}{\linewidth}
\noindent{\scriptsize\textsf{82: multiplication}}
\begin{align*}
r_{82}={}&-z_{78}z_{97}+z_{98}
\end{align*}
\end{minipage}\par
\par\noindent\begin{minipage}{\linewidth}
\noindent{\scriptsize\textsf{83: multiplication}}
\begin{align*}
r_{83}={}&-z_{88}z_{98}+z_{99}
\end{align*}
\end{minipage}\par
\par\noindent\begin{minipage}{\linewidth}
\noindent{\scriptsize\textsf{84: multiplication}}
\begin{align*}
r_{84}={}&-z_{90}z_{99}+z_{100}
\end{align*}
\end{minipage}\par
\par\noindent\begin{minipage}{\linewidth}
\noindent{\scriptsize\textsf{85: multiplication}}
\begin{align*}
r_{85}={}&-z_{100}^{2}+z_{101}
\end{align*}
\end{minipage}\par
\par\noindent\begin{minipage}{\linewidth}
\noindent{\scriptsize\textsf{86: multiplication}}
\begin{align*}
r_{86}={}&-1-z_{2}-z_{2}z_{18}+z_{102}
\end{align*}
\end{minipage}\par
\par\noindent\begin{minipage}{\linewidth}
\noindent{\scriptsize\textsf{87: multiplication}}
\begin{align*}
r_{87}={}&-1-z_{16}z_{102}+z_{103}
\end{align*}
\end{minipage}\par
\par\noindent\begin{minipage}{\linewidth}
\noindent{\scriptsize\textsf{88: nonnegative\_square\_difference}}
\begin{align*}
r_{88}={}&-2z_{1}^{2}+2z_{1}z_{14}+2z_{1}z_{15}+2z_{1}z_{103}+2z_{1}z_{104}-z_{14}^{2}-2z_{14}z_{103}-z_{15}^{2}-2z_{15}z_{104}\\[-2pt]
&+z_{19}z_{106}-z_{81}z_{101}-z_{83}z_{101}-z_{103}^{2}-z_{104}^{2}
\end{align*}
\end{minipage}\par
\par\noindent\begin{minipage}{\linewidth}
\noindent{\scriptsize\textsf{89: multiplication}}
\begin{align*}
r_{89}={}&-1-z_{2}z_{17}-z_{17}z_{102}+z_{104}
\end{align*}
\end{minipage}\par
\par\noindent\begin{minipage}{\linewidth}
\noindent{\scriptsize\textsf{90: addition}}
\begin{align*}
r_{90}={}&-9z_{10}-9z_{11}-4z_{21}-2z_{24}-4z_{26}-2z_{29}-2z_{31}-4z_{32}-2z_{34}-2z_{36}-4z_{39}-2z_{40}-2z_{43}-2z_{47}\\[-2pt]
&-4z_{68}-2z_{69}+z_{105}
\end{align*}
\end{minipage}\par
\par\noindent\begin{minipage}{\linewidth}
\noindent{\scriptsize\textsf{91: multiplication}}
\begin{align*}
r_{91}={}&-z_{107}z_{110}+z_{110}+z_{112}
\end{align*}
\end{minipage}\par
\par\noindent\begin{minipage}{\linewidth}
\noindent{\scriptsize\textsf{92: multiplication}}
\begin{align*}
r_{92}={}&-z_{5}z_{109}+z_{108}-z_{109}z_{112}-z_{111}
\end{align*}
\end{minipage}\par
\par\noindent\begin{minipage}{\linewidth}
\noindent{\scriptsize\textsf{93: multiplication}}
\begin{align*}
r_{93}={}&-z_{19}z_{114}+z_{117}
\end{align*}
\end{minipage}\par
\par\noindent\begin{minipage}{\linewidth}
\noindent{\scriptsize\textsf{94: multiplication}}
\begin{align*}
r_{94}={}&z_{113}-z_{115}z_{117}-z_{116}
\end{align*}
\end{minipage}\par
\par\noindent\begin{minipage}{\linewidth}
\noindent{\scriptsize\textsf{95: multiplication}}
\begin{align*}
r_{95}={}&-z_{117}z_{119}+z_{118}-z_{120}
\end{align*}
\end{minipage}\par
\par\noindent\begin{minipage}{\linewidth}
\noindent{\scriptsize\textsf{96: multiplication}}
\begin{align*}
r_{96}={}&-z_{117}z_{122}+z_{121}-z_{123}
\end{align*}
\end{minipage}\par
\par\noindent\begin{minipage}{\linewidth}
\noindent{\scriptsize\textsf{97: multiplication}}
\begin{align*}
r_{97}={}&-z_{117}z_{125}+z_{124}-z_{126}
\end{align*}
\end{minipage}\par
\par\noindent\begin{minipage}{\linewidth}
\noindent{\scriptsize\textsf{98: multiplication}}
\begin{align*}
r_{98}={}&-z_{19}z_{128}+z_{131}
\end{align*}
\end{minipage}\par
\par\noindent\begin{minipage}{\linewidth}
\noindent{\scriptsize\textsf{99: multiplication}}
\begin{align*}
r_{99}={}&z_{127}-z_{129}z_{131}-z_{130}
\end{align*}
\end{minipage}\par
\par\noindent\begin{minipage}{\linewidth}
\noindent{\scriptsize\textsf{100: multiplication}}
\begin{align*}
r_{100}={}&-z_{131}z_{133}+z_{132}-z_{134}
\end{align*}
\end{minipage}\par
\par\noindent\begin{minipage}{\linewidth}
\noindent{\scriptsize\textsf{101: multiplication}}
\begin{align*}
r_{101}={}&-z_{19}z_{136}+z_{139}
\end{align*}
\end{minipage}\par
\par\noindent\begin{minipage}{\linewidth}
\noindent{\scriptsize\textsf{102: multiplication}}
\begin{align*}
r_{102}={}&z_{135}-z_{137}z_{139}-z_{138}
\end{align*}
\end{minipage}\par
\par\noindent\begin{minipage}{\linewidth}
\noindent{\scriptsize\textsf{103: multiplication}}
\begin{align*}
r_{103}={}&-z_{139}z_{141}+z_{140}-z_{142}
\end{align*}
\end{minipage}\par
\par\noindent\begin{minipage}{\linewidth}
\noindent{\scriptsize\textsf{104: multiplication}}
\begin{align*}
r_{104}={}&-z_{19}z_{109}+z_{146}
\end{align*}
\end{minipage}\par
\par\noindent\begin{minipage}{\linewidth}
\noindent{\scriptsize\textsf{105: multiplication}}
\begin{align*}
r_{105}={}&z_{143}-z_{144}z_{146}-z_{145}
\end{align*}
\end{minipage}\par
\par\noindent\begin{minipage}{\linewidth}
\noindent{\scriptsize\textsf{106: multiplication}}
\begin{align*}
r_{106}={}&-z_{117}z_{148}+z_{147}-z_{149}
\end{align*}
\end{minipage}\par
\par\noindent\begin{minipage}{\linewidth}
\noindent{\scriptsize\textsf{107: multiplication}}
\begin{align*}
r_{107}={}&z_{150}-2z_{150}z_{152}+z_{153}
\end{align*}
\end{minipage}\par
\par\noindent\begin{minipage}{\linewidth}
\noindent{\scriptsize\textsf{108: multiplication}}
\begin{align*}
r_{108}={}&-z_{150}z_{153}-z_{153}+z_{154}
\end{align*}
\end{minipage}\par
\par\noindent\begin{minipage}{\linewidth}
\noindent{\scriptsize\textsf{109: shared.last\_norm}}
\begin{align*}
r_{109}={}&2z_{151}+z_{151}^{2}+z_{153}^{2}-z_{154}^{2}
\end{align*}
\end{minipage}\par
\par\noindent\begin{minipage}{\linewidth}
\noindent{\scriptsize\textsf{110: multiplication}}
\begin{align*}
r_{110}={}&1+z_{151}-z_{151}z_{156}-z_{156}+z_{157}
\end{align*}
\end{minipage}\par
\par\noindent\begin{minipage}{\linewidth}
\noindent{\scriptsize\textsf{111: index0.norm1}}
\begin{align*}
r_{111}={}&z_{155}^{2}-2z_{156}+z_{156}^{2}-z_{157}^{2}
\end{align*}
\end{minipage}\par
\par\noindent\begin{minipage}{\linewidth}
\noindent{\scriptsize\textsf{112: multiplication}}
\begin{align*}
r_{112}={}&z_{4}^{2}-2z_{4}z_{151}-2z_{151}+z_{158}
\end{align*}
\end{minipage}\par
\par\noindent\begin{minipage}{\linewidth}
\noindent{\scriptsize\textsf{113: multiplication}}
\begin{align*}
r_{113}={}&1+z_{4}-z_{4}z_{156}+z_{6}-z_{155}-z_{156}+z_{157}+z_{158}z_{159}
\end{align*}
\end{minipage}\par
\par\noindent\begin{minipage}{\linewidth}
\noindent{\scriptsize\textsf{114: multiplication}}
\begin{align*}
r_{114}={}&-z_{151}z_{161}-z_{161}+z_{162}
\end{align*}
\end{minipage}\par
\par\noindent\begin{minipage}{\linewidth}
\noindent{\scriptsize\textsf{115: index1.norm1}}
\begin{align*}
r_{115}={}&-1+z_{160}^{2}+z_{161}^{2}-z_{162}^{2}
\end{align*}
\end{minipage}\par
\par\noindent\begin{minipage}{\linewidth}
\noindent{\scriptsize\textsf{116: multiplication}}
\begin{align*}
r_{116}={}&1-2z_{5}+z_{5}^{2}-2z_{5}z_{151}+z_{163}
\end{align*}
\end{minipage}\par
\par\noindent\begin{minipage}{\linewidth}
\noindent{\scriptsize\textsf{117: multiplication}}
\begin{align*}
r_{117}={}&-z_{5}z_{161}+z_{19}-z_{160}+z_{162}+z_{163}z_{164}
\end{align*}
\end{minipage}\par
\par\noindent\begin{minipage}{\linewidth}
\noindent{\scriptsize\textsf{118: base2.T\_definition}}
\begin{align*}
r_{118}={}&-1+2z_{151}z_{224}-z_{165}+2z_{224}-z_{224}^{2}
\end{align*}
\end{minipage}\par
\par\noindent\begin{minipage}{\linewidth}
\noindent{\scriptsize\textsf{119: multiplication}}
\begin{align*}
r_{119}={}&-1+z_{107}-z_{160}-z_{161}z_{224}+z_{162}+z_{165}z_{166}
\end{align*}
\end{minipage}\par
\par\noindent\begin{minipage}{\linewidth}
\noindent{\scriptsize\textsf{120: multiplication}}
\begin{align*}
r_{120}={}&1+z_{151}-z_{151}z_{168}-z_{168}+z_{169}
\end{align*}
\end{minipage}\par
\par\noindent\begin{minipage}{\linewidth}
\noindent{\scriptsize\textsf{121: index2.norm1}}
\begin{align*}
r_{121}={}&z_{167}^{2}-2z_{168}+z_{168}^{2}-z_{169}^{2}
\end{align*}
\end{minipage}\par
\par\noindent\begin{minipage}{\linewidth}
\noindent{\scriptsize\textsf{122: multiplication}}
\begin{align*}
r_{122}={}&1-2z_{107}+z_{107}^{2}-2z_{107}z_{151}+z_{170}
\end{align*}
\end{minipage}\par
\par\noindent\begin{minipage}{\linewidth}
\noindent{\scriptsize\textsf{123: multiplication}}
\begin{align*}
r_{123}={}&z_{107}-z_{107}z_{168}+z_{108}-z_{167}+z_{169}+z_{170}z_{171}
\end{align*}
\end{minipage}\par
\par\noindent\begin{minipage}{\linewidth}
\noindent{\scriptsize\textsf{124: multiplication}}
\begin{align*}
r_{124}={}&4-4z_{107}+z_{107}^{2}-2z_{107}z_{151}+2z_{151}+z_{172}
\end{align*}
\end{minipage}\par
\par\noindent\begin{minipage}{\linewidth}
\noindent{\scriptsize\textsf{125: multiplication}}
\begin{align*}
r_{125}={}&-1+z_{107}-z_{107}z_{156}+z_{109}-z_{155}+z_{156}+z_{157}+z_{172}z_{173}
\end{align*}
\end{minipage}\par
\par\noindent\begin{minipage}{\linewidth}
\noindent{\scriptsize\textsf{126: multiplication}}
\begin{align*}
r_{126}={}&1+z_{151}-z_{151}z_{175}-z_{175}+z_{176}
\end{align*}
\end{minipage}\par
\par\noindent\begin{minipage}{\linewidth}
\noindent{\scriptsize\textsf{127: index3.norm1}}
\begin{align*}
r_{127}={}&z_{174}^{2}-2z_{175}+z_{175}^{2}-z_{176}^{2}
\end{align*}
\end{minipage}\par
\par\noindent\begin{minipage}{\linewidth}
\noindent{\scriptsize\textsf{128: multiplication}}
\begin{align*}
r_{128}={}&z_{107}-z_{107}z_{175}+z_{113}+z_{170}z_{177}-z_{174}+z_{176}
\end{align*}
\end{minipage}\par
\par\noindent\begin{minipage}{\linewidth}
\noindent{\scriptsize\textsf{129: multiplication}}
\begin{align*}
r_{129}={}&-2z_{151}-z_{151}^{2}+z_{178}
\end{align*}
\end{minipage}\par
\par\noindent\begin{minipage}{\linewidth}
\noindent{\scriptsize\textsf{130: multiplication}}
\begin{align*}
r_{130}={}&-z_{107}z_{155}-z_{107}z_{157}+z_{114}+z_{151}z_{157}+z_{155}-z_{156}z_{178}+2z_{157}+z_{172}z_{179}+z_{178}
\end{align*}
\end{minipage}\par
\par\noindent\begin{minipage}{\linewidth}
\noindent{\scriptsize\textsf{131: multiplication}}
\begin{align*}
r_{131}={}&-z_{151}z_{181}-z_{181}+z_{182}
\end{align*}
\end{minipage}\par
\par\noindent\begin{minipage}{\linewidth}
\noindent{\scriptsize\textsf{132: index5.norm1}}
\begin{align*}
r_{132}={}&-1+z_{180}^{2}+z_{181}^{2}-z_{182}^{2}
\end{align*}
\end{minipage}\par
\par\noindent\begin{minipage}{\linewidth}
\noindent{\scriptsize\textsf{133: multiplication}}
\begin{align*}
r_{133}={}&-z_{107}z_{181}+z_{118}+z_{170}z_{183}-z_{180}+z_{182}
\end{align*}
\end{minipage}\par
\par\noindent\begin{minipage}{\linewidth}
\noindent{\scriptsize\textsf{134: multiplication}}
\begin{align*}
r_{134}={}&-z_{151}z_{185}-z_{185}+z_{186}
\end{align*}
\end{minipage}\par
\par\noindent\begin{minipage}{\linewidth}
\noindent{\scriptsize\textsf{135: index6.norm1}}
\begin{align*}
r_{135}={}&-1+z_{184}^{2}+z_{185}^{2}-z_{186}^{2}
\end{align*}
\end{minipage}\par
\par\noindent\begin{minipage}{\linewidth}
\noindent{\scriptsize\textsf{136: multiplication}}
\begin{align*}
r_{136}={}&-z_{107}z_{185}+z_{121}+z_{170}z_{187}-z_{184}+z_{186}
\end{align*}
\end{minipage}\par
\par\noindent\begin{minipage}{\linewidth}
\noindent{\scriptsize\textsf{137: multiplication}}
\begin{align*}
r_{137}={}&-z_{151}z_{189}-z_{189}+z_{190}
\end{align*}
\end{minipage}\par
\par\noindent\begin{minipage}{\linewidth}
\noindent{\scriptsize\textsf{138: index7.norm1}}
\begin{align*}
r_{138}={}&-1+z_{188}^{2}+z_{189}^{2}-z_{190}^{2}
\end{align*}
\end{minipage}\par
\par\noindent\begin{minipage}{\linewidth}
\noindent{\scriptsize\textsf{139: multiplication}}
\begin{align*}
r_{139}={}&-z_{107}z_{189}+z_{124}+z_{170}z_{191}-z_{188}+z_{190}
\end{align*}
\end{minipage}\par
\par\noindent\begin{minipage}{\linewidth}
\noindent{\scriptsize\textsf{140: multiplication}}
\begin{align*}
r_{140}={}&-z_{151}z_{193}-z_{193}+z_{194}
\end{align*}
\end{minipage}\par
\par\noindent\begin{minipage}{\linewidth}
\noindent{\scriptsize\textsf{141: index8.norm1}}
\begin{align*}
r_{141}={}&-1+z_{192}^{2}+z_{193}^{2}-z_{194}^{2}
\end{align*}
\end{minipage}\par
\par\noindent\begin{minipage}{\linewidth}
\noindent{\scriptsize\textsf{142: multiplication}}
\begin{align*}
r_{142}={}&-z_{107}z_{193}+z_{127}+z_{170}z_{195}-z_{192}+z_{194}
\end{align*}
\end{minipage}\par
\par\noindent\begin{minipage}{\linewidth}
\noindent{\scriptsize\textsf{143: multiplication}}
\begin{align*}
r_{143}={}&-1-z_{151}-z_{151}z_{197}-z_{197}+z_{198}
\end{align*}
\end{minipage}\par
\par\noindent\begin{minipage}{\linewidth}
\noindent{\scriptsize\textsf{144: index9.norm1}}
\begin{align*}
r_{144}={}&z_{196}^{2}+2z_{197}+z_{197}^{2}-z_{198}^{2}
\end{align*}
\end{minipage}\par
\par\noindent\begin{minipage}{\linewidth}
\noindent{\scriptsize\textsf{145: multiplication}}
\begin{align*}
r_{145}={}&1-z_{107}-z_{107}z_{197}+z_{128}+z_{172}z_{199}-z_{196}+z_{197}+z_{198}
\end{align*}
\end{minipage}\par
\par\noindent\begin{minipage}{\linewidth}
\noindent{\scriptsize\textsf{146: multiplication}}
\begin{align*}
r_{146}={}&-z_{151}z_{201}-z_{201}+z_{202}
\end{align*}
\end{minipage}\par
\par\noindent\begin{minipage}{\linewidth}
\noindent{\scriptsize\textsf{147: index10.norm1}}
\begin{align*}
r_{147}={}&-1+z_{200}^{2}+z_{201}^{2}-z_{202}^{2}
\end{align*}
\end{minipage}\par
\par\noindent\begin{minipage}{\linewidth}
\noindent{\scriptsize\textsf{148: multiplication}}
\begin{align*}
r_{148}={}&-z_{107}z_{201}+z_{132}+z_{170}z_{203}-z_{200}+z_{202}
\end{align*}
\end{minipage}\par
\par\noindent\begin{minipage}{\linewidth}
\noindent{\scriptsize\textsf{149: multiplication}}
\begin{align*}
r_{149}={}&-z_{151}z_{205}-z_{205}+z_{206}
\end{align*}
\end{minipage}\par
\par\noindent\begin{minipage}{\linewidth}
\noindent{\scriptsize\textsf{150: index11.norm1}}
\begin{align*}
r_{150}={}&-1+z_{204}^{2}+z_{205}^{2}-z_{206}^{2}
\end{align*}
\end{minipage}\par
\par\noindent\begin{minipage}{\linewidth}
\noindent{\scriptsize\textsf{151: multiplication}}
\begin{align*}
r_{151}={}&-z_{107}z_{205}+z_{135}+z_{170}z_{207}-z_{204}+z_{206}
\end{align*}
\end{minipage}\par
\par\noindent\begin{minipage}{\linewidth}
\noindent{\scriptsize\textsf{152: multiplication}}
\begin{align*}
r_{152}={}&-z_{151}z_{209}-z_{209}+z_{210}
\end{align*}
\end{minipage}\par
\par\noindent\begin{minipage}{\linewidth}
\noindent{\scriptsize\textsf{153: index12.norm1}}
\begin{align*}
r_{153}={}&-1+z_{208}^{2}+z_{209}^{2}-z_{210}^{2}
\end{align*}
\end{minipage}\par
\par\noindent\begin{minipage}{\linewidth}
\noindent{\scriptsize\textsf{154: multiplication}}
\begin{align*}
r_{154}={}&-z_{107}z_{209}+z_{136}+z_{172}z_{211}-z_{208}+z_{209}+z_{210}
\end{align*}
\end{minipage}\par
\par\noindent\begin{minipage}{\linewidth}
\noindent{\scriptsize\textsf{155: multiplication}}
\begin{align*}
r_{155}={}&-z_{151}z_{213}-z_{213}+z_{214}
\end{align*}
\end{minipage}\par
\par\noindent\begin{minipage}{\linewidth}
\noindent{\scriptsize\textsf{156: index13.norm1}}
\begin{align*}
r_{156}={}&-1+z_{212}^{2}+z_{213}^{2}-z_{214}^{2}
\end{align*}
\end{minipage}\par
\par\noindent\begin{minipage}{\linewidth}
\noindent{\scriptsize\textsf{157: multiplication}}
\begin{align*}
r_{157}={}&-z_{107}z_{213}+z_{140}+z_{170}z_{215}-z_{212}+z_{214}
\end{align*}
\end{minipage}\par
\par\noindent\begin{minipage}{\linewidth}
\noindent{\scriptsize\textsf{158: multiplication}}
\begin{align*}
r_{158}={}&-z_{151}z_{217}-z_{217}+z_{218}
\end{align*}
\end{minipage}\par
\par\noindent\begin{minipage}{\linewidth}
\noindent{\scriptsize\textsf{159: index14.norm1}}
\begin{align*}
r_{159}={}&-1+z_{216}^{2}+z_{217}^{2}-z_{218}^{2}
\end{align*}
\end{minipage}\par
\par\noindent\begin{minipage}{\linewidth}
\noindent{\scriptsize\textsf{160: multiplication}}
\begin{align*}
r_{160}={}&-z_{107}z_{217}+z_{143}+z_{170}z_{219}-z_{216}+z_{218}
\end{align*}
\end{minipage}\par
\par\noindent\begin{minipage}{\linewidth}
\noindent{\scriptsize\textsf{161: multiplication}}
\begin{align*}
r_{161}={}&-z_{151}z_{221}-z_{221}+z_{222}
\end{align*}
\end{minipage}\par
\par\noindent\begin{minipage}{\linewidth}
\noindent{\scriptsize\textsf{162: index15.norm1}}
\begin{align*}
r_{162}={}&-1+z_{220}^{2}+z_{221}^{2}-z_{222}^{2}
\end{align*}
\end{minipage}\par
\par\noindent\begin{minipage}{\linewidth}
\noindent{\scriptsize\textsf{163: multiplication}}
\begin{align*}
r_{163}={}&-z_{107}z_{221}+z_{147}+z_{170}z_{223}-z_{220}+z_{222}
\end{align*}
\end{minipage}\par
\par\noindent\begin{minipage}{\linewidth}
\noindent{\scriptsize\textsf{164: shared.frame\_definition}}
\begin{align*}
r_{164}={}&-8+4z_{0}-7z_{0}^{2}-8z_{0}z_{1}-8z_{0}z_{2}-8z_{0}z_{3}-2z_{0}z_{4}-5z_{1}^{2}-8z_{1}z_{2}-8z_{1}z_{3}-2z_{2}\\[-2pt]
&-5z_{2}^{2}-8z_{2}z_{3}-4z_{3}^{2}-2z_{4}^{2}-z_{5}^{2}-z_{6}^{2}+2z_{10}-z_{10}^{2}-z_{11}^{2}-z_{12}^{2}-z_{13}^{2}\\[-2pt]
&-z_{14}^{2}-z_{15}^{2}-z_{16}^{2}-z_{17}^{2}-z_{18}^{2}-z_{19}^{2}-z_{105}^{2}+2z_{107}-2z_{107}^{2}-z_{108}^{2}\\[-2pt]
&-z_{109}^{2}-z_{113}^{2}-z_{114}^{2}-z_{118}^{2}-z_{121}^{2}-z_{124}^{2}-z_{127}^{2}-z_{128}^{2}-z_{132}^{2}\\[-2pt]
&-z_{135}^{2}-z_{136}^{2}-z_{140}^{2}-z_{143}^{2}-z_{147}^{2}+z_{150}-z_{224}^{2}
\end{align*}
\end{minipage}\par
\par\noindent\begin{minipage}{\linewidth}
\noindent{\scriptsize\textsf{165: self\_radix.base\_definition}}
\begin{align*}
r_{165}={}&-z_{5}z_{10}-z_{5}z_{11}-z_{5}z_{12}-z_{5}z_{13}-z_{5}z_{14}-z_{5}z_{15}-z_{5}z_{16}-z_{5}z_{17}-z_{5}z_{18}+z_{224}
\end{align*}
\end{minipage}\par
\par\noindent\begin{minipage}{\linewidth}
\noindent{\scriptsize\textsf{166: cofinal.horizon}}
\begin{align*}
r_{166}={}&-17+z_{0}-z_{225}^{2}
\end{align*}
\end{minipage}\par
\par\noindent\begin{minipage}{\linewidth}
\noindent{\scriptsize\textsf{167: cofinal.beta\_base}}
\begin{align*}
r_{167}={}&z_{2}+z_{6}-z_{6}^{2}
\end{align*}
\end{minipage}\par
\par\noindent\begin{minipage}{\linewidth}
\noindent{\scriptsize\textsf{168: cap.main\_Ay}}
\begin{align*}
r_{168}={}&-z_{151}z_{227}-z_{227}+z_{238}
\end{align*}
\end{minipage}\par
\par\noindent\begin{minipage}{\linewidth}
\noindent{\scriptsize\textsf{169: cap.main\_norm}}
\begin{align*}
r_{169}={}&-1+z_{226}^{2}+z_{227}^{2}-z_{238}^{2}
\end{align*}
\end{minipage}\par
\par\noindent\begin{minipage}{\linewidth}
\noindent{\scriptsize\textsf{170: cap.common\_Ay}}
\begin{align*}
r_{170}={}&-z_{151}z_{229}-z_{229}+z_{239}
\end{align*}
\end{minipage}\par
\par\noindent\begin{minipage}{\linewidth}
\noindent{\scriptsize\textsf{171: cap.common\_norm}}
\begin{align*}
r_{171}={}&-1+z_{228}^{2}+z_{229}^{2}-z_{239}^{2}
\end{align*}
\end{minipage}\par
\par\noindent\begin{minipage}{\linewidth}
\noindent{\scriptsize\textsf{172: cap.BR}}
\begin{align*}
r_{172}={}&-z_{231}-z_{231}z_{232}+z_{240}
\end{align*}
\end{minipage}\par
\par\noindent\begin{minipage}{\linewidth}
\noindent{\scriptsize\textsf{173: cap.comparison\_norm}}
\begin{align*}
r_{173}={}&-1+z_{230}^{2}+z_{231}^{2}-z_{240}^{2}
\end{align*}
\end{minipage}\par
\par\noindent\begin{minipage}{\linewidth}
\noindent{\scriptsize\textsf{174: cap.Y2}}
\begin{align*}
r_{174}={}&-z_{227}^{2}+z_{241}
\end{align*}
\end{minipage}\par
\par\noindent\begin{minipage}{\linewidth}
\noindent{\scriptsize\textsf{175: cap.common\_multiple}}
\begin{align*}
r_{175}={}&z_{229}-z_{234}z_{241}-z_{241}
\end{align*}
\end{minipage}\par
\par\noindent\begin{minipage}{\linewidth}
\noindent{\scriptsize\textsf{176: cap.comparison\_constructor}}
\begin{align*}
r_{176}={}&-4z_{229}-8z_{229}z_{233}+z_{232}
\end{align*}
\end{minipage}\par
\par\noindent\begin{minipage}{\linewidth}
\noindent{\scriptsize\textsf{177: cap.comparison\_modulus}}
\begin{align*}
r_{177}={}&-z_{151}-z_{228}z_{235}+z_{232}
\end{align*}
\end{minipage}\par
\par\noindent\begin{minipage}{\linewidth}
\noindent{\scriptsize\textsf{178: cap.abscissa\_residue}}
\begin{align*}
r_{178}={}&-z_{226}-z_{228}z_{236}+z_{230}
\end{align*}
\end{minipage}\par
\par\noindent\begin{minipage}{\linewidth}
\noindent{\scriptsize\textsf{179: cap.index\_residue}}
\begin{align*}
r_{179}={}&-z_{150}-4z_{227}z_{237}+z_{231}
\end{align*}
\end{minipage}\par
\par\noindent\begin{minipage}{\linewidth}
\noindent{\scriptsize\textsf{180: cap.pooled\_index\_residue}}
\begin{align*}
r_{180}={}&9+6z_{0}^{2}+8z_{0}z_{1}+8z_{0}z_{2}+8z_{0}z_{3}+2z_{0}z_{4}-2z_{0}z_{156}-4z_{0}z_{161}-2z_{0}z_{168}+5z_{1}^{2}\\[-2pt]
&+8z_{1}z_{2}+8z_{1}z_{3}-4z_{1}z_{161}-2z_{1}z_{209}+5z_{2}^{2}+8z_{2}z_{3}-4z_{2}z_{161}-2z_{2}z_{197}+4z_{3}^{2}\\[-2pt]
&-4z_{3}z_{161}+z_{4}^{2}-2z_{4}z_{168}+z_{10}^{2}-2z_{10}z_{175}+z_{11}^{2}-2z_{11}z_{181}+z_{12}^{2}-2z_{12}z_{185}\\[-2pt]
&+z_{13}^{2}-2z_{13}z_{189}+z_{14}^{2}-2z_{14}z_{193}+z_{15}^{2}-2z_{15}z_{201}+z_{16}^{2}-2z_{16}z_{205}+z_{17}^{2}\\[-2pt]
&-2z_{17}z_{213}+z_{18}^{2}-2z_{18}z_{217}+z_{105}^{2}-2z_{105}z_{221}-z_{151}z_{242}+z_{156}^{2}+z_{161}^{2}\\[-2pt]
&+z_{168}^{2}-z_{174}+z_{175}^{2}-z_{180}+z_{181}^{2}-z_{184}+z_{185}^{2}-z_{188}+z_{189}^{2}-z_{192}+z_{193}^{2}\\[-2pt]
&+z_{197}^{2}-z_{200}+z_{201}^{2}-z_{204}+z_{205}^{2}+z_{209}^{2}-z_{212}+z_{213}^{2}-z_{216}+z_{217}^{2}+z_{221}^{2}
\end{align*}
\end{minipage}\par
\par\noindent\begin{minipage}{\linewidth}
\noindent{\scriptsize\textsf{181: domain\_guard.0}}
\begin{align*}
r_{181}={}&-1+z_{1}-z_{243}^{2}-z_{244}^{2}-z_{245}-z_{245}^{2}
\end{align*}
\end{minipage}\par
\par\noindent\begin{minipage}{\linewidth}
\noindent{\scriptsize\textsf{182: domain\_guard.1}}
\begin{align*}
r_{182}={}&-1-z_{3}z_{5}+z_{5}+z_{6}-z_{246}^{2}-z_{247}^{2}-z_{248}-z_{248}^{2}
\end{align*}
\end{minipage}\par
\par\noindent\begin{minipage}{\linewidth}
\noindent{\scriptsize\textsf{183: domain\_guard.2}}
\begin{align*}
r_{183}={}&z_{3}z_{5}-z_{6}-z_{249}^{2}-z_{250}^{2}-z_{251}-z_{251}^{2}
\end{align*}
\end{minipage}\par
\par\noindent\begin{minipage}{\linewidth}
\noindent{\scriptsize\textsf{184: domain\_guard.12}}
\begin{align*}
r_{184}={}&-2-z_{5}+z_{107}-z_{252}^{2}-z_{253}^{2}-z_{254}-z_{254}^{2}
\end{align*}
\end{minipage}\par
\par\noindent\begin{minipage}{\linewidth}
\noindent{\scriptsize\textsf{185: domain\_guard.13}}
\begin{align*}
r_{185}={}&1+z_{151}-z_{255}^{2}-z_{256}^{2}-z_{257}-z_{257}^{2}
\end{align*}
\end{minipage}\par
\par\noindent\begin{minipage}{\linewidth}
\noindent{\scriptsize\textsf{186: domain\_guard.14}}
\begin{align*}
r_{186}={}&z_{109}z_{111}-z_{111}-z_{111}^{2}-z_{258}^{2}-z_{259}^{2}-z_{260}-z_{260}^{2}
\end{align*}
\end{minipage}\par
\par\noindent\begin{minipage}{\linewidth}
\noindent{\scriptsize\textsf{187: domain\_guard.15}}
\begin{align*}
r_{187}={}&z_{114}z_{116}-z_{116}-z_{116}^{2}-z_{261}^{2}-z_{262}^{2}-z_{263}-z_{263}^{2}
\end{align*}
\end{minipage}\par
\par\noindent\begin{minipage}{\linewidth}
\noindent{\scriptsize\textsf{188: domain\_guard.16}}
\begin{align*}
r_{188}={}&z_{114}z_{120}-z_{120}-z_{120}^{2}-z_{264}^{2}-z_{265}^{2}-z_{266}-z_{266}^{2}
\end{align*}
\end{minipage}\par
\par\noindent\begin{minipage}{\linewidth}
\noindent{\scriptsize\textsf{189: domain\_guard.17}}
\begin{align*}
r_{189}={}&z_{114}z_{123}-z_{123}-z_{123}^{2}-z_{267}^{2}-z_{268}^{2}-z_{269}-z_{269}^{2}
\end{align*}
\end{minipage}\par
\par\noindent\begin{minipage}{\linewidth}
\noindent{\scriptsize\textsf{190: domain\_guard.18}}
\begin{align*}
r_{190}={}&z_{114}z_{126}-z_{126}-z_{126}^{2}-z_{270}^{2}-z_{271}^{2}-z_{272}-z_{272}^{2}
\end{align*}
\end{minipage}\par
\par\noindent\begin{minipage}{\linewidth}
\noindent{\scriptsize\textsf{191: domain\_guard.19}}
\begin{align*}
r_{191}={}&z_{128}z_{130}-z_{130}-z_{130}^{2}-z_{273}^{2}-z_{274}^{2}-z_{275}-z_{275}^{2}
\end{align*}
\end{minipage}\par
\par\noindent\begin{minipage}{\linewidth}
\noindent{\scriptsize\textsf{192: domain\_guard.20}}
\begin{align*}
r_{192}={}&z_{128}z_{134}-z_{134}-z_{134}^{2}-z_{276}^{2}-z_{277}^{2}-z_{278}-z_{278}^{2}
\end{align*}
\end{minipage}\par
\par\noindent\begin{minipage}{\linewidth}
\noindent{\scriptsize\textsf{193: domain\_guard.21}}
\begin{align*}
r_{193}={}&z_{136}z_{138}-z_{138}-z_{138}^{2}-z_{279}^{2}-z_{280}^{2}-z_{281}-z_{281}^{2}
\end{align*}
\end{minipage}\par
\par\noindent\begin{minipage}{\linewidth}
\noindent{\scriptsize\textsf{194: domain\_guard.22}}
\begin{align*}
r_{194}={}&z_{136}z_{142}-z_{142}-z_{142}^{2}-z_{282}^{2}-z_{283}^{2}-z_{284}-z_{284}^{2}
\end{align*}
\end{minipage}\par
\par\noindent\begin{minipage}{\linewidth}
\noindent{\scriptsize\textsf{195: domain\_guard.23}}
\begin{align*}
r_{195}={}&z_{109}z_{145}-z_{145}-z_{145}^{2}-z_{285}^{2}-z_{286}^{2}-z_{287}-z_{287}^{2}
\end{align*}
\end{minipage}\par
\par\noindent\begin{minipage}{\linewidth}
\noindent{\scriptsize\textsf{196: domain\_guard.24}}
\begin{align*}
r_{196}={}&z_{114}z_{149}-z_{149}-z_{149}^{2}-z_{288}^{2}-z_{289}^{2}-z_{290}-z_{290}^{2}
\end{align*}
\end{minipage}\par
\par\noindent\begin{minipage}{\linewidth}
\noindent{\scriptsize\textsf{197: domain\_guard.25}}
\begin{align*}
r_{197}={}&z_{234}-z_{291}^{2}-z_{292}^{2}-z_{293}-z_{293}^{2}
\end{align*}
\end{minipage}\par
\par\noindent\begin{minipage}{\linewidth}
\noindent{\scriptsize\textsf{198: domain\_guard.26}}
\begin{align*}
r_{198}={}&-5-z_{150}+2z_{156}-z_{156}^{2}-z_{161}^{2}+2z_{168}-z_{168}^{2}+2z_{175}-z_{175}^{2}-z_{181}^{2}-z_{185}^{2}\\[-2pt]
&-z_{189}^{2}-z_{193}^{2}-2z_{197}-z_{197}^{2}-z_{201}^{2}-z_{205}^{2}-z_{209}^{2}-z_{213}^{2}-z_{217}^{2}-z_{221}^{2}\\[-2pt]
&+z_{227}-z_{294}^{2}-z_{295}^{2}-z_{296}-z_{296}^{2}
\end{align*}
\end{minipage}\par
\par\noindent\begin{minipage}{\linewidth}
\noindent{\scriptsize\textsf{199: orientation.product\_2}}
\begin{align*}
r_{199}={}&-z_{10}z_{11}+z_{11}+z_{297}
\end{align*}
\end{minipage}\par
\par\noindent\begin{minipage}{\linewidth}
\noindent{\scriptsize\textsf{200: orientation.product\_3}}
\begin{align*}
r_{200}={}&-z_{12}z_{297}+z_{298}
\end{align*}
\end{minipage}\par
\par\noindent\begin{minipage}{\linewidth}
\noindent{\scriptsize\textsf{201: orientation.product\_4}}
\begin{align*}
r_{201}={}&-z_{13}z_{298}+z_{299}
\end{align*}
\end{minipage}\par
\par\noindent\begin{minipage}{\linewidth}
\noindent{\scriptsize\textsf{202: orientation.product\_5}}
\begin{align*}
r_{202}={}&-z_{14}z_{299}+z_{300}
\end{align*}
\end{minipage}\par
\par\noindent\begin{minipage}{\linewidth}
\noindent{\scriptsize\textsf{203: orientation.product\_6}}
\begin{align*}
r_{203}={}&-z_{15}z_{300}+z_{301}
\end{align*}
\end{minipage}\par
\par\noindent\begin{minipage}{\linewidth}
\noindent{\scriptsize\textsf{204: orientation.product\_7}}
\begin{align*}
r_{204}={}&-z_{16}z_{301}+z_{302}
\end{align*}
\end{minipage}\par
\par\noindent\begin{minipage}{\linewidth}
\noindent{\scriptsize\textsf{205: orientation.fused\_strict\_product}}
\begin{align*}
r_{205}={}&-1+z_{17}z_{302}-z_{303}^{2}-z_{304}^{2}-z_{305}-z_{305}^{2}
\end{align*}
\end{minipage}\par
\endgroup


\section{Source identities and comparison inventory}
The comparison and retrieval pin is \srcpin. The direct parent comes from
\path{131-shared-pell-caps-round16/direct/integer324/}; the matched graph parent comes from
\path{127-graph-certificates/graph/principal/}. The 306-coordinate source is newly constructed here.
The complete per-file inventory is \path{MANIFEST.sha256}; this appendix records the normative equation identities.
\begingroup\small
\par\noindent\path{direct/parent/integer324-equation.json}\
\mbox{36,840 bytes.} SHA-256: \path{32618060c22765f54423b53f9faef437d736494da82aef309efe96210d6cbadd}.\par\smallskip
\par\noindent\path{direct/integer306/equation.json}\
\mbox{35,585 bytes.} SHA-256: \path{d51d75d2c6398fa3410b19ae43ba8d80f86c668ad670a803f6777a2af221c189}.\par\smallskip
\par\noindent\path{direct/integer306/expanded_polynomial.json}\
\mbox{114,236 bytes.} SHA-256: \path{c3b4b6c7646ba6a92e02aa24fc9d64e62d642e47046a6d38c1e01df7357d293a}.\par\smallskip
\par\noindent\path{graph/matched1252/polynomial.json}\
\mbox{656,104 bytes.} SHA-256: \path{66a29766c5b1107c36674c38215b02afaba1802606cb4f1a0584cc24302c2bf8}.\par\smallskip
\par\noindent\path{graph/principal985/polynomial.json}\
\mbox{1,984,728 bytes.} SHA-256: \path{33e35bb2f0354cbcc47bac27056c8fb560c508cb78ce04b1bb1dff9adcf7a635}.\par\smallskip
\par\noindent\path{graph/sparse1122/polynomial.json}\
\mbox{580,440 bytes.} SHA-256: \path{69edd5df14b21dee3960da9cb86f5723e0e99aaf13814922cda965493ab77d9f}.\par\smallskip
\endgroup
The retained Lean interface files identify the inherited logical calibration. They are reference sources, not a complete standalone Lean project; their full transitive imports were not rebuilt for this release.


\clearpage
\begin{thebibliography}{9}
\bibitem{graph127} ProveIt research packet 127, \emph{Shared arithmetic
certificates and graph energies for Con(ZFC)}, 24 October 2026.
\url{https://github.com/VladimirReshetnikov/ProveIt/tree/1c5367452a0bd4dd7e9e4bceb1508b62dbd2ff61/Computability/DiophantineConsistency/con-zfc-diophantine-forms/127-graph-certificates}.
\bibitem{cap131} ProveIt research packet 131, \emph{Shared Pell caps and
coprime product equations for ZFC consistency}, retained at the same pin.
\url{https://github.com/VladimirReshetnikov/ProveIt/tree/1c5367452a0bd4dd7e9e4bceb1508b62dbd2ff61/Computability/DiophantineConsistency/con-zfc-diophantine-forms/131-shared-pell-caps-round16}.
\bibitem{repooverview} ProveIt, \emph{Con(ZFC) in Diophantine form},
catalogue and active-research review, snapshot \srcpin.
\url{https://github.com/VladimirReshetnikov/ProveIt/tree/1c5367452a0bd4dd7e9e4bceb1508b62dbd2ff61/Computability/DiophantineConsistency}.
\bibitem{jones} J. P. Jones, \emph{Universal Diophantine equation},
Journal of Symbolic Logic 47 (1982), 549--571.
\url{https://doi.org/10.2307/2273588}.
\bibitem{bayer} J. Bayer, M. David, M. Hassler, Y. Matiyasevich, and
D. Schleicher, \emph{Diophantine Equations over $\mathbb Z$: Universal
Bounds and Parallel Formalization}, arXiv:2506.20909v3, 22 October 2026.
\url{https://arxiv.org/abs/2506.20909v3}.
\end{thebibliography}
\end{document}
